% ============================================================
%  Lucrare de Licență – Universitatea Politehnica Timișoara
%  Student: Craiu Bogdan
%  Coordonator: Ș.l. dr. ing. Cătălin Iapă
% ============================================================
\documentclass[12pt, a4paper, twoside]{report}

% ---------- pachete esențiale ----------
\usepackage{lmodern}
\usepackage[T1]{fontenc}
\usepackage[utf8]{inputenc}
\usepackage[romanian]{babel}

% ---------- caractere române\c{s}ti: comma-below (U+0218-021B) lipse\c{s}te din T1;
%            map\u{a}m la cedil\u{a} care e suportat\u{a} nativ de T1+lmodern ----------
\DeclareUnicodeCharacter{0218}{\c{S}}
\DeclareUnicodeCharacter{0219}{\c{s}}
\DeclareUnicodeCharacter{021A}{\c{T}}
\DeclareUnicodeCharacter{021B}{\c{t}}
\usepackage{geometry}
\usepackage{graphicx}
\usepackage{amsmath, amssymb, amsfonts}
\usepackage{listings}
\usepackage{xcolor}
\usepackage{hyperref}
\usepackage{booktabs}
\usepackage{array}
\usepackage{tabularx}
\usepackage{longtable}
\usepackage{float}
\usepackage{caption}
\usepackage{subcaption}
\usepackage{fancyhdr}
\usepackage{setspace}
\usepackage{titlesec}
\usepackage{enumitem}
\usepackage{microtype}
\usepackage{tocloft}
\usepackage{appendix}
\usepackage{indentfirst}
\usepackage{pdfpages}

% ---------- geometrie pagină ----------
\geometry{
  a4paper,
  top=2.5cm, bottom=2.5cm,
  left=3cm, right=2.5cm,
  headheight=2.5cm
}

% ---------- spațiere ----------
\setstretch{1.15}

% ---------- cale imagini ----------
\graphicspath{{images/}}

% ---------- culori cod ----------
\definecolor{codegreen}{rgb}{0,0.6,0}
\definecolor{codegray}{rgb}{0.5,0.5,0.5}
\definecolor{codepurple}{rgb}{0.58,0,0.82}
\definecolor{backcolour}{rgb}{0.97,0.97,0.97}
\definecolor{keywordblue}{rgb}{0.13,0.29,0.53}

\lstdefinestyle{mystyle}{
  backgroundcolor=\color{backcolour},
  commentstyle=\color{codegreen}\itshape,
  keywordstyle=\color{keywordblue}\bfseries,
  numberstyle=\tiny\color{codegray},
  stringstyle=\color{codepurple},
  basicstyle=\ttfamily\footnotesize,
  breakatwhitespace=false,
  breaklines=true,
  captionpos=b,
  keepspaces=true,
  numbers=left,
  numbersep=5pt,
  showspaces=false,
  showstringspaces=false,
  showtabs=false,
  tabsize=2,
  frame=single,
  rulecolor=\color{codegray}
}
\lstset{style=mystyle}

% ---------- stil pagini ----------
\pagestyle{fancy}
\fancyhf{}
\renewcommand{\chaptermark}[1]{\markboth{\thechapter.\ #1}{}}
\fancyhead[L]{\scriptsize\parbox[t]{0.65\headwidth}{%
  Universitatea Politehnica Timișoara \\
  Calculatoare și Tehnologia Informației, 2026 \\
  Craiu Bogdan \\
  Aplicație Mobilă RPG pentru Studiu, Concentrare și Jocuri Inteligente cu AI}}
\fancyhead[R]{\raisebox{-1.0cm}{
\includegraphics[height=1.6cm]{logo_UPT}}}
\fancyfoot[C]{\thepage}
\renewcommand{\headrulewidth}{0.4pt}

\fancypagestyle{plain}{%
  \fancyhf{}
  \fancyhead[L]{\scriptsize\parbox[t]{0.65\headwidth}{%
    Universitatea Politehnica Timișoara \\
    Calculatoare și Tehnologia Informației, 2026 \\
    Craiu Bogdan \\
    Aplicație Mobilă RPG pentru Studiu, Concentrare și Jocuri Inteligente cu AI}}
  \fancyhead[R]{\raisebox{-1.0cm}{
\includegraphics[height=1.6cm]{logo_UPT}}}
  \fancyfoot[C]{\thepage}
  \renewcommand{\headrulewidth}{0.4pt}
}

% ---------- hyperref ----------
\hypersetup{
  colorlinks=true,
  linkcolor=blue!60!black,
  citecolor=green!50!black,
  urlcolor=blue!70!black,
  pdfauthor={Craiu Bogdan},
  pdftitle={Aplicație Mobilă RPG pentru Studiu și Concentrare},
}

% ---------- font ----------
\renewcommand{\rmdefault}{phv}

% ---------- localizare tabele ----------
\AtBeginDocument{%
  \renewcommand{\tablename}{Tabelul}
}

% ---------- numerotare ecuații ----------
\numberwithin{equation}{section}

% ---------- titluri capitole ----------
\titleformat{\chapter}[display]
  {\normalfont\LARGE\bfseries}{\chaptertitlename\ \thechapter}{18pt}{\LARGE}
\titlespacing{\chapter}{0pt}{50pt}{30pt}

% ============================================================
\begin{document}

% ============================================================
%  PAGINA DE TITLU
% ============================================================
\begin{titlepage}
\thispagestyle{empty}

\noindent
\begin{minipage}[t]{0.60\textwidth}
  {\large\bfseries Calculatoare și Tehnologia Informației}\\[0.3cm]
  {\large\bfseries 2026}
\end{minipage}%
\hfill
\begin{minipage}[t]{0.35\textwidth}
  \begin{flushright}
    
\includegraphics[width=\linewidth]{logo_UPT}
  \end{flushright}
\end{minipage}

\vspace{5cm}
\begin{center}
  {\bfseries\LARGE Aplicație Mobilă RPG pentru Studiu,\\[0.5cm]
  Concentrare și Jocuri Inteligente cu AI}\\[1cm]
  {\bfseries\Large --- FocusQuest ---}
\end{center}

\vspace{4cm}

\noindent{\large\bfseries Candidat: Craiu Bogdan}\\[0.5cm]
\noindent{\large\bfseries Coordonator științific: Ș.l. dr. ing. Cătălin Iapă}

\vfill

\begin{center}
  {\large Sesiunea: Iulie 2026}
\end{center}
\setcounter{page}{1}
\end{titlepage}

% ============================================================
%  REZUMAT (RO)
% ============================================================
\chapter*{Rezumat}
\addcontentsline{toc}{chapter}{Rezumat}
\markboth{Rezumat}{Rezumat}

Prezenta lucrare de licență descrie proiectarea, implementarea și testarea unei aplicații mobile
cross-platform denumită \textbf{FocusQuest}, dezvoltată cu React Native și TypeScript.
Aplicația îmbină elemente de joc de rol (RPG) cu instrumente de productivitate, fiind
destinată oricui este interesat de productivitate, studiu și jocuri inteligente cu AI, mai
ales elevilor care urmează să susțină admiterea pentru intrarea la Universitatea Politehnica Timișoara.

Principalele contribuții ale lucrării sunt:
\begin{itemize}[leftmargin=1.5cm]
  \item Integrarea unui model de rețea neuronală convoluțională (\textit{ResNet-50 INT8}) rulat
    \textit{on-device} prin ONNX Runtime React Native, utilizat pentru clasificarea imaginilor
    în jurnalul digital al utilizatorului.
  \item Integrarea motorului de șah \textit{Maia Chess} — un model antrenat să imite stilul
    de joc uman la diverse niveluri Elo (1100–1500) — rulat complet local, fără conexiune la
    internet.
  \item Un sistem de economie virtuală bazat pe monede câștigate prin sesiuni de concentrare,
    jocuri și activități de învățare, utilizat pentru deblocarea conținutului de dificultate
    mai ridicată.
  \item Un modul de învățare structurat pe trei discipline matematice (algebră, analiză şi
    trigonometrie), cu probleme matematice şi soluții pas-cu-pas
    randate cu KaTeX.
  \item Persistența stării aplicației folosind AsyncStorage și un gestionar global de stare
    bazat pe React Context API cu \texttt{useReducer}.
  \item Un sistem de gamificare de tip RPG ce cuprinde experiență (XP) cu niveluri progresive,
    quest-uri zilnice cu resetare automată și efecte sonore retro sintetizate, proiectat
    pentru a susține motivația pe termen lung.
  \item Un modul de examen simulat (12 întrebări extrase aleatoriu, câte 4 din fiecare
    disciplină) și un set de puzzle-uri de șah de tip ``mat în $N$ mutări'' validate
    algoritmic cu \texttt{chess.js}, pentru consolidarea cunoștințelor și a gândirii tactice.
\end{itemize}

Aplicația a fost testată pe dispozitive Android și respectă arhitectura componentelor React
Native, cu o separare clară a logicii de afaceri față de stratul de prezentare.

\textbf{Cuvinte cheie:} React Native, TypeScript, ONNX Runtime, ResNet-50, Maia Chess,
gamificare, niveluri XP, examen simulat, puzzle-uri de șah, aplicație mobilă, productivitate,
algoritmi de backtracking, KaTeX.

\newpage

% ============================================================
%  ABSTRACT (EN)
% ============================================================
\chapter*{Abstract}
\addcontentsline{toc}{chapter}{Abstract}
\markboth{Abstract}{Abstract}

This bachelor thesis describes the design, implementation and testing of a cross-platform
mobile application called \textbf{FocusQuest}, built with React Native and TypeScript.
The application combines role-playing game (RPG) mechanics with productivity tools, aimed at
anyone interested in productivity, learning and intelligent games.

The main contributions are:
\begin{itemize}[leftmargin=1.5cm]
  \item Integration of a convolutional neural network (\textit{ResNet-50 INT8}) running
    \textit{on-device} via ONNX Runtime React Native, used to classify images in the user's
    digital journal.
  \item Integration of the \textit{Maia Chess} engine — a model trained to mimic human chess
    playing style at various Elo levels (1100–1500) — running fully locally, without an
    internet connection.
  \item A virtual economy system based on coins earned through focus sessions, games and
    learning activities, used to unlock higher-difficulty content.
  \item A learning module structured around three mathematical disciplines (algebra, analysis
    and trigonometry), with mathematical problems and step-by-step
    solutions rendered with KaTeX.
  \item Application state persistence using AsyncStorage and a global state manager based on
    the React Context API with \texttt{useReducer}.
  \item An RPG-style gamification system comprising experience points (XP) with progressive
    levels, auto-resetting daily quests and synthesized retro sound effects, designed to
    sustain long-term motivation.
  \item A mock-exam module (12 randomly drawn questions, four per subject) and a set of
    ``mate-in-$N$'' chess puzzles algorithmically validated with \texttt{chess.js}, for
    reinforcing knowledge and tactical thinking.
\end{itemize}

\textbf{Keywords:} React Native, TypeScript, ONNX Runtime, ResNet-50, Maia Chess,
gamification, XP levels, mock exam, chess puzzles, mobile application, productivity,
backtracking algorithms, KaTeX.

\newpage

% ============================================================
%  CUPRINS
% ============================================================
\tableofcontents
\newpage

% ============================================================
%  LISTA FIGURILOR / TABELE
% ============================================================
\listoftables
\addcontentsline{toc}{chapter}{Lista tabelelor}
\listoffigures
\addcontentsline{toc}{chapter}{Lista figurilor}

% ============================================================
%  CAPITOLUL 1 – INTRODUCERE
% ============================================================
\chapter{Introducere}

\section{Motivație și context}

Accesul tinerilor la instrumente moderne de studiu s-a diversificat semnificativ în ultimul
deceniu. Telefoanele inteligente sunt omniprezente în rândul elevilor și studenților, iar
aplicațiile mobile dedicate productivității și educației reprezintă un segment în creștere
rapidă al pieței de software. Cu toate acestea, există o lipsă acută de aplicații care să
combine în mod integrat:
\begin{enumerate}
  \item Tehnici de productivitate validate din punct de vedere \textit{\c{s}tiin\c{t}ific} (ex. tehnica Pomodoro);
  \item Conținut de pregătire pentru examene specifice sistemului de învățământ românesc;
  \item Mecanisme de gamificare care să mențină motivația pe termen lung;
  \item Modele de inteligență artificială care să funcționeze \textit{offline}, pe dispozitiv.
\end{enumerate}

\textbf{FocusQuest} răspunde acestei nevoi prin crearea unui ecosistem mobil coerent în care
studiul și distracția coexistă. Aplicația este construită cu un design vizual retro-neon RPG,
conceput să fie atractiv pentru generația Z, obișnuită cu estetica jocurilor video indie.

\section{Obiectivele lucrării}

Obiectivele principale ale acestei lucrări de licență sunt:
\begin{itemize}
  \item Proiectarea și implementarea unei aplicații mobile React Native complete, cu navigare
    multi-ecran, gestionare globală a stării și persistență locală a datelor.
  \item Integrarea a două modele ONNX \textit{on-device}: ResNet-50 pentru clasificarea
    imaginilor și Maia Chess pentru jocul de șah cu inteligență artificială.
  \item Construirea unui sistem de gamificare funcțional cu monede, experiență (XP) și
    niveluri, quest-uri zilnice, realizări deblocabile, inventar și dificultăți progresive.
  \item Furnizarea unui modul educațional cu exerciții matematice structurate, complet
    cu răspunsuri multiple şi soluții detaliate randate cu KaTeX, precum și a unui mod de
    examen simulat și a unui set de puzzle-uri de șah tactice.
  \item Demonstrarea viabilității rulării modelelor de rețele neuronale direct pe dispozitive
    mobile cu resurse limitate.
\end{itemize}

\section{Structura lucrării}

Lucrarea este organizată în nouă capitole:
\begin{itemize}
  \item \textbf{Capitolul 2} prezintă stadiul actual al domeniului (State of the Art),
    analizând aplicații similare existente și comparând funcționalitățile acestora cu
    FocusQuest.
  \item \textbf{Capitolul 3} descrie tehnologiile și instrumentele utilizate în dezvoltare.
  \item \textbf{Capitolul 4} prezintă arhitectura generală a aplicației.
  \item \textbf{Capitolul 5} detaliază sistemul de gestiune a stării și persistența datelor.
  \item \textbf{Capitolul 6} descrie modulele funcționale principale ale aplicației.
  \item \textbf{Capitolul 7} analizează integrarea modelelor de inteligență artificială.
  \item \textbf{Capitolul 8} prezintă testarea și evaluarea aplicației.
  \item \textbf{Capitolul 9} concluzionează lucrarea și propune direcții viitoare de
    dezvoltare.
\end{itemize}

% ============================================================
%  CAPITOLUL 2 – STATE OF THE ART
% ============================================================
\chapter{Stadiul actual al domeniului (State of the Art)}

\section{Aplicații de productivitate și concentrare}

\subsection{Forest – Stay Focused}

Forest~\cite{forest} este una dintre cele mai cunoscute aplicații de productivitate cu element de
gamificare. Utilizatorul plantează un copac virtual care crește cât timp nu folosește
telefonul; dacă iese din aplicație, copacul moare. Aplicația colaborează cu organizații
de plantare a copacilor reali, ceea ce adaugă o dimensiune de impact social.

\textbf{Comparație cu FocusQuest:} Forest se concentrează exclusiv pe tehnica Pomodoro
și nu integrează niciun conținut educațional, niciun joc inteligent și niciun model AI.
FocusQuest extinde conceptul de concentrare cu un ecosistem complet: sesiunile de focus
generează monede care permit accesul la jocuri și conținut educațional.

\subsection{Todoist și Notion}

Todoist și Notion sunt instrumente de gestionare a sarcinilor și note-taking, utilizate
frecvent de studenți pentru organizarea studiului. Deși extrem de versatile, aceste aplicații
nu oferă gamificare intrinsecă, nu integrează modele AI locale şi nu sunt orientate spre
practică structurată de exerciții matematice.

\subsection{Flora – Green Focus}

Flora urmează un model similar cu Forest, cu accent pe productivitate în grup și provocări
sociale. Nu include elemente RPG, nu are jocuri integrate și nu are funcționalitate offline
bazată pe AI.

\section{Aplicații de învățare și pregătire pentru examene}

\subsection{Khan Academy}

Khan Academy oferă un curriculum vast, inclusiv matematică avansată, cu exerciții
interactive și progres vizibil. Platforma este gratuită și bine structurată, dar:
\begin{itemize}
  \item Nu este orientată spre conținut specific admiterii la universitățile românești;
  \item Nu are elemente de gamificare RPG;
  \item Nu integrează jocuri sau un sistem de monede;
  \item Necesită conexiune la internet.
\end{itemize}

\subsection{Brilliant.org}

Brilliant.org oferă cursuri interactive de matematică și informatică cu un nivel ridicat de
interactivitate. Interfața este modernă și angajantă, dar platforma este orientată spre
gândire conceptuală, nu spre pregătirea pentru examinări de tip grilă specifice.

\subsection{Duolingo}

Duolingo~\cite{duolingo} este exemplul canonic de aplicație educațională cu gamificare: streak-uri zilnice,
monede virtuale, nivele și personaje animate. FocusQuest se inspiră din principiile de design
ale Duolingo (streak-uri, monede, recompense), adaptate contextului unui RPG și extins cu
funcționalități AI locale.

\textbf{Diferența cheie față de FocusQuest:} FocusQuest combină zona de practică de exerciții
cu zona de jocuri (șah AI, Sudoku) și cu zona de productivitate (timer focus, geolocation)
într-o singură aplicație cu o economie virtuală unificată.

\section{Aplicații de șah cu AI}

\subsection{Chess.com și Lichess}

Chess.com și Lichess~\cite{lichess} sunt platformele dominante de șah online. Ambele oferă oponenți AI
la multiple niveluri de dificultate. Lichess este open-source și oferă motorul Stockfish~\cite{stockfish},
unul dintre cei mai puternici motori de șah din lume.

\textbf{Diferența față de FocusQuest:} FocusQuest nu utilizează Stockfish (care joacă la
nivel de campion mondial). În schimb, integrează \textbf{Maia Chess}, un model antrenat să
joace \textit{uman}, comițând greșeli tipice ale jucătorilor la niveluri Elo între 1100 și 1500.
Aceasta face experiența mult mai potrivită pentru jucători începători și intermediari.
Mai mult, FocusQuest rulează modelul complet offline, fără server.

\subsection{Magnus Trainer}

Magnus Trainer este o aplicație de antrenament la șah creată de campionul mondial Magnus
Carlsen. Oferă lecții, puzzle-uri și jocuri contra AI. Aplicația este premium și nu
integrează module de studiu sau productivitate.

\section{Aplicații cu clasificare AI a imaginilor}

\subsection{Google Lens}

Google Lens oferă recunoaștere de imagini de înaltă calitate, rulând pe serverele Google.
Necesită conexiune la internet și nu poate fi integrată offline într-o aplicație mobilă.

\subsection{Microsoft Azure Custom Vision}

Serviciu cloud pentru clasificare de imagini custom. Necesită abonament și conexiune la
internet; nu este potrivit pentru scenarii complet offline.

\textbf{Diferența față de FocusQuest:} FocusQuest rulează ResNet-50 INT8 direct pe
dispozitivul utilizatorului prin ONNX Runtime, fără niciun server extern, asigurând
confidențialitatea datelor și funcționalitate offline completă.

\section{Aplicații RPG educaționale}

\subsection{Habitica}

Habitica~\cite{habitica} transformă gestionarea obiceiurilor și sarcinilor zilnice într-un RPG: utilizatorul
are un personaj, câștigă experiență și obiecte, luptă cu monștri în grupuri. Este cea mai
apropiată aplicație conceptual de FocusQuest, dar:
\begin{itemize}
  \item Nu integrează modele AI locale;
  \item Nu oferă jocuri inteligente (șah, Sudoku);
  \item Nu are conținut educațional structurat cu exerciții specifice;
  \item Nu are un timer de concentrare dedicat cu animații de personaj.
\end{itemize}

\subsection{ClassDojo}

ClassDojo este orientat spre sala de clasă și gestionat de profesori, nu de elevi în mod
autonom. Nu are elemente de jocuri inteligente sau AI.

\section{Rezumat comparativ}

\begin{longtable}{>{\bfseries}p{3.5cm} p{1.4cm} p{1.4cm} p{1.4cm} p{1.4cm} p{1.4cm} p{1.4cm}}
\caption{Comparație funcționalități -- FocusQuest vs. aplicații similare}
\label{tab:comparison}\\
\toprule
Funcționalitate & Focus Quest & Forest & Habitica & Duolingo & Chess.com & Lichess \\
\midrule
\endfirsthead
\toprule
Funcționalitate & Focus Quest & Forest & Habitica & Duolingo & Chess.com & Lichess \\
\midrule
\endhead
\bottomrule
\endfoot
Timer focus & Da & Da & Parțial & Nu & Nu & Nu \\
Gamificare RPG & Da & Parțial & Da & Parțial & Nu & Nu \\
Conținut edu. & Da & Nu & Nu & Da & Nu & Nu \\
Joc șah AI & Da & Nu & Nu & Nu & Da & Da \\
AI local (offline) & Da & Nu & Nu & Nu & Nu & Nu \\
Sudoku & Da & Nu & Nu & Nu & Nu & Nu \\
Jurnal cu AI & Da & Nu & Nu & Nu & Nu & Nu \\
Geolocation & Da & Nu & Nu & Nu & Nu & Nu \\
Muzică ambientală & Da & Nu & Nu & Nu & Nu & Nu \\
Inventar & Da & Nu & Da & Nu & Nu & Nu \\
Realizări & Da & Parțial & Da & Da & Da & Da \\
Funcționare offline & Da & Da & Nu & Parțial & Nu & Parțial \\
\bottomrule
\end{longtable}

Din tabelul \ref{tab:comparison} reiese că FocusQuest ocupă o nișă unică prin combinarea
tuturor acestor funcționalități într-o singură aplicație. Nicio altă aplicație existentă
nu reunește timere de concentrare, jocuri inteligente cu AI local, conținut educațional
specific și un sistem RPG de gamificare.

% ============================================================
%  CAPITOLUL 3 – TEHNOLOGII UTILIZATE
% ============================================================
\chapter{Tehnologii și instrumente utilizate}

\section{React Native}

React Native~\cite{reactnative} este un framework open-source dezvoltat de Meta (Facebook) care permite
crearea de aplicații mobile native pentru iOS și Android folosind JavaScript și TypeScript.
Spre deosebire de abordările hibride (ex. Ionic, Cordova), React Native compilează
componentele la UI nativ al platformei, nu la WebView, rezultând o performanță apropiată
de aplicațiile native.

Versiunea utilizată în proiect este \textbf{React Native 0.82.0} împreună cu
\textbf{React 19.1.1}.

Avantajele care au determinat alegerea React Native:
\begin{itemize}
  \item Baza de cod partajată pentru iOS și Android;
  \item Ecosistem bogat de biblioteci (navigare, stocare, imagini, sunet);
  \item Suport nativ pentru TypeScript;
  \item Compatibilitate cu ONNX Runtime prin pachetul \texttt{onnxruntime-react-native}.
\end{itemize}

\section{TypeScript}

TypeScript~\cite{typescript} este un superset al JavaScript care adaugă tipizare statică. Întregul proiect
este scris în TypeScript, ceea ce oferă:
\begin{itemize}
  \item Detectarea timpurie a erorilor la compilare;
  \item Autocompletare și refactorizare mai bune în IDE;
  \item Documentare implicită a structurilor de date complexe (ex. \texttt{GameState}).
\end{itemize}

\section{React Navigation}

Navigarea în aplicație este gestionată de biblioteca \textbf{React Navigation 7.x}~\cite{reactnavigation},
utilizând două tipuri de navigatori:
\begin{itemize}
  \item \texttt{createBottomTabNavigator} — pentru cele cinci ecrane principale
    (Focus, Arcade, Learning, Journal, Profile);
  \item \texttt{createNativeStackNavigator} — pentru ecranele secundare
    (Achievements, Statistics, Inventory, ExerciseList, ExerciseDetail).
\end{itemize}

\section{AsyncStorage}

\textbf{@react-native-async-storage/async-storage}~\cite{asyncstorage} (versiunea 2.1.0) este biblioteca
de stocare cheie-valoare persistentă utilizată pentru salvarea stării aplicației între
sesiuni. Datele sunt serializate în format JSON și stocate local pe dispozitiv sub cheia
\texttt{retro-arcade-study-state-v1}.

\section{ONNX Runtime React Native}

\textbf{onnxruntime-react-native}~\cite{onnxruntime} (versiunea 1.20.0) este un pachet care expune
ONNX Runtime~\cite{onnxformat} — un motor de inferență pentru modele de machine learning în format ONNX
— aplicațiilor React Native. Permite rularea modelelor complet local, fără server.

Funcționalitățile utilizate:
\begin{itemize}
  \item \texttt{InferenceSession.create()} — încărcarea modelului ONNX din fișier;
  \item \texttt{session.run(feeds)} — executarea inferenței cu un dicționar de tensori
    de intrare;
  \item \texttt{Tensor} — crearea tensorilor de tip \texttt{float32} și \texttt{int64}.
\end{itemize}

\section{chess.js}

\textbf{chess.js}~\cite{chessjs} (versiunea 1.0.0) este o bibliotecă JavaScript de logică a jocului
de șah. Oferă:
\begin{itemize}
  \item Validarea mutărilor legale;
  \item Detecția situațiilor de șah, mat, pat;
  \item Conversia pozițiilor în format FEN (Forsyth-Edwards Notation);
  \item Generarea listei complete de mutări legale în format UCI.
\end{itemize}

\section{react-native-sound}

\textbf{react-native-sound} (versiunea 0.13.0) este un modul pentru redarea fișierelor
audio locale. Este utilizat pentru serviciul de muzică de fundal (muzică lo-fi) integrat
în ecranul principal. Fișierele audio sunt stocate în directorul
\texttt{android/app/src/main/res/raw/}.

\section{react-native-image-picker}

\textbf{react-native-image-picker} (versiunea 7.2.1) permite accesarea galeriei foto
sau a camerei dispozitivului. Este utilizat în modulul jurnal (FlashcardsScreen) pentru
adăugarea imaginilor la intrările de jurnal.

\section{react-native-fs}

\textbf{react-native-fs} (versiunea 2.20.0) oferă acces la sistemul de fișiere al
dispozitivului. Este utilizat pentru copierea fișierelor de model ONNX din bundlul
de resurse Android în directorul de cache, de unde pot fi accesate de ONNX Runtime.

\section{react-native-webview}

\textbf{react-native-webview} (versiunea 13.16.1) este utilizat de componenta
\texttt{MathRenderer} pentru randarea formulelor matematice LaTeX cu biblioteca KaTeX,
încărcată în interiorul WebView-ului dintr-o rețea de distribuție de conținut (CDN).

\section{KaTeX}

\textbf{KaTeX}~\cite{katex} (versiunea 0.16.11) este o bibliotecă JavaScript de randare a formulelor
matematice LaTeX, utilizată în ecranele de exerciții pentru afișarea problemelor și
soluțiilor cu formatare matematică corectă. Biblioteca este încărcată într-un WebView embed
direct de pe CDN-ul public \texttt{jsdelivr} (fișierele \texttt{katex.min.css},
\texttt{katex.min.js} și extensia \texttt{auto-render}). Randarea formulelor matematice
reprezintă, prin urmare, \textbf{singura funcționalitate a aplicației care necesită o
conexiune activă la internet}; toate celelalte module — inclusiv cele două modele de
inteligență artificială — funcționează complet \textit{offline} (detalii în
Secțiunea~\ref{sec:online-dependency}).

\section{Alte pachete}

Dependențele suplimentare utilizate în proiect sunt prezentate în Tabelul~\ref{tab:deps}.

\begin{table}[H]
\centering
\caption{Dependențe suplimentare ale proiectului}
\label{tab:deps}
\begin{tabular}{lll}
\toprule
\textbf{Pachet} & \textbf{Versiune} & \textbf{Scop} \\
\midrule
react-native-gesture-handler & 2.16.2 & Gesture recognition \\
react-native-reanimated & 4.1.3 & Animații UI \\
react-native-safe-area-context & 5.5.2 & Zone sigure pe ecran \\
react-native-screens & 4.4.0 & Ecrane native optimizate \\
@react-native-community/geolocation & 3.4.0 & Geolocalizare GPS \\
jpeg-js & 0.4.4 & Decodificare JPEG în JavaScript \\
buffer & 6.0.3 & Manipulare buffere binare \\
\bottomrule
\end{tabular}
\end{table}

% ============================================================
%  CAPITOLUL 4 – ARHITECTURA APLICAȚIEI
% ============================================================
\chapter{Arhitectura aplicației}

\section{Prezentare generală}

FocusQuest urmează o arhitectură pe straturi, specifică aplicațiilor React Native moderne:

\begin{enumerate}
  \item \textbf{Stratul de prezentare (UI)} — componente React (TSX) ce definesc interfața;
  \item \textbf{Stratul de logică} — hook-uri custom, context și reducer;
  \item \textbf{Stratul de date} — AsyncStorage pentru persistență și ONNX Runtime pentru
    inferență AI;
  \item \textbf{Stratul de navigare} — React Navigation cu combinare de navigatori.
\end{enumerate}

\section{Structura directorului de proiect}

\begin{lstlisting}[language=bash, caption={Structura principală a proiectului}]
FocusQuest/
  App.tsx                   # Punct de intrare principal
  package.json              # Dependente si scripturi
  tsconfig.json             # Configurare TypeScript
  android/                  # Cod nativ Android
  ios/                      # Cod nativ iOS
  src/
    assets/
      skeletonAnimations.ts # Definitii animatii sprite
    components/             # Componente reutilizabile UI
    constants/              # Date statice si cataloage
    context/
      GameContext.tsx        # Store global (Context + Reducer)
    hooks/                  # Hook-uri custom React
    navigation/
      RootNavigator.tsx     # Configurare navigare
    screens/                # Ecrane principale
    services/
      MusicPlayer.ts        # Serviciu muzica de fundal
      SoundEffects.ts       # Serviciu efecte sonore (chiptune)
    theme/
      colors.ts             # Paleta de culori
    types/
      achievements.ts       # Definitii realizari
\end{lstlisting}

\section{Sistemul de navigare}

Aplicația utilizează o arhitectură de navigare pe două niveluri:

\begin{itemize}
  \item \textbf{Stack Navigator (rădăcină)} — conține TabNavigator ca ecran principal
    și ecranele secundare (Achievements, Statistics, Inventory, ExerciseList, ExerciseDetail,
    Quiz, Puzzles) ca ecrane de stivă.
  \item \textbf{Tab Navigator (principal)} — cinci tab-uri bottom: Focus, Arcade, Learn,
    Journal, Profile.
\end{itemize}

Tipurile de parametri ale navigatoarelor sunt definite explicit în TypeScript:

\begin{lstlisting}[language=Java, caption={Tipuri parametri navigatoare}]
export type RootTabParamList = {
  Home: undefined;
  Arcade: undefined;
  Learning: undefined;
  Flashcards: undefined;
  Profile: undefined;
};

export type RootStackParamList = {
  MainTabs: undefined;
  Achievements: undefined;
  Statistics: undefined;
  Inventory: undefined;
  ExerciseList: { subject: LearningSubject };
  ExerciseDetail: { subject: LearningSubject; exerciseId: string };
  Quiz: undefined;
  Puzzles: undefined;
};
\end{lstlisting}

\section{Fluxul de date al aplicației}

Fluxul de date respectă principiul unidirecțional specific arhitecturii Flux/Redux:

\begin{enumerate}
  \item Utilizatorul interacționează cu un component UI (ex. apasă butonul de Start Timer);
  \item Componentul apelează o funcție din contextul global (ex. \texttt{completeSession});
  \item Funcția dispatch-uiește un \texttt{Action} către reducer;
  \item Reducer-ul calculează noua stare pe baza acțiunii și a stării curente;
  \item Noua stare este salvată în AsyncStorage;
  \item Toate componentele care consumă contextul sunt re-randate cu noua stare.
\end{enumerate}

\section{Tema vizuală}

Paleta de culori a aplicației este definită în \texttt{src/theme/colors.ts} și respectă
o estetică retro-neon RPG (Tabelul~\ref{tab:colors}):

\begin{table}[H]
\centering
\caption{Paleta de culori a aplicației}
\label{tab:colors}
\begin{tabular}{lll}
\toprule
\textbf{Variabilă} & \textbf{Cod hex} & \textbf{Utilizare} \\
\midrule
\texttt{midnight} & \texttt{\#0b1026} & Fundal principal \\
\texttt{neonBlue} & \texttt{\#38bdf8} & Accente informative \\
\texttt{neonPink} & \texttt{\#f472b6} & Accent primar, indicator activ \\
\texttt{neonYellow} & \texttt{\#facc15} & Monede, titluri, evidențieri \\
\texttt{neonGreen} & \texttt{\#34d399} & Acțiuni pozitive (adăugare) \\
\texttt{electricPurple} & \texttt{\#a855f7} & Accentuări secundare \\
\texttt{softWhite} & \texttt{\#f9fafb} & Text principal \\
\texttt{danger} & \texttt{\#f87171} & Erori, avertizări \\
\bottomrule
\end{tabular}
\end{table}

\section{Reutilizarea componentelor: tabla de șah partajată}

Cele două moduri de șah ale aplicației — \textbf{Cosmic Chess Gauntlet} (cu motor minimax)
și \textbf{Maia AI Chess Challenge} (cu rețea neurală) — partajează exact aceeași
reprezentare vizuală a tablei. Pentru a respecta principiul \textit{DRY} (\textit{Don't
Repeat Yourself}) și a evita duplicarea codului, randarea tablei a fost extrasă într-o
componentă de prezentare reutilizabilă, \texttt{ChessBoard} (\texttt{src/components/ChessBoard.tsx}).

Componenta este complet \textit{stateless}: primește prin proprietăți starea tablei
(\texttt{boardState}), pătratul selectat, lista de mutări legale și un \textit{callback}
de interacțiune, fiind responsabilă exclusiv de desenarea grilei $8\times8$, a etichetelor
de rânduri/coloane, a evidențierii pătratului selectat și a indicatorilor de mutare legală:

\begin{lstlisting}[language=Java, caption={Interfața componentei reutilizabile ChessBoard}]
type Props = {
  boardState: (Piece | null)[][];
  selectedSquare: string | null;
  legalTargets: string[];
  onSquarePress: (square: string) => void;
  orientation?: 'white' | 'black'; // permite rotirea tablei
};
\end{lstlisting}

Fiecare arenă de șah își păstrează exclusiv propria logică (motorul AI, dificultățile,
recompensele), delegând în întregime afișarea tablei către \texttt{ChessBoard}. Refactorizarea
a eliminat aproximativ 330 de linii de cod duplicat, concentrând randarea într-un singur loc
ușor de întreținut. Proprietatea \texttt{orientation} permite, suplimentar, rotirea tablei
pentru ca jucătorul să poată juca și cu piesele negre (Secțiunea~\ref{sec:play-black}).

\section{Inițializarea aplicației (\texttt{App.tsx})}

Punctul de intrare al aplicației este \texttt{App.tsx}, care construiește ierarhia
de provideri React și gestionează starea de hidrare:

\begin{lstlisting}[language=Java, caption={Ierarhia de provideri în App.tsx}]
export default function App() {
  return (
    <GestureHandlerRootView style={styles.root}>
      <SafeAreaProvider>
        <StatusBar barStyle="light-content" />
        <GameProvider>
          <AppShell />
        </GameProvider>
      </SafeAreaProvider>
    </GestureHandlerRootView>
  );
}
\end{lstlisting}

Fiecare container din ierarhie îndeplinește un rol distinct:
\begin{itemize}
  \item \texttt{GestureHandlerRootView} \textemdash{} necesar pentru
    \texttt{react-native-gesture-handler}, care trebuie să fie primul element
    din arborele de componente;
  \item \texttt{SafeAreaProvider} \textemdash{} furnizează insets-urile pentru
    zona sigură (notch, home indicator) tuturor ecranelor descendente;
  \item \texttt{StatusBar barStyle="light-content"} \textemdash{} menține
    textul barei de stare alb, compatibil cu fundalul întunecat al aplicației;
  \item \texttt{GameProvider} \textemdash{} inițializează starea globală și
    restaurează datele din AsyncStorage.
\end{itemize}

Componenta internă \texttt{AppShell} implementează o \textbf{gardă de hidrare}:
atât timp cât \texttt{isHydrated = false} (GameProvider nu a terminat de restaurat
starea din AsyncStorage), se afișează un spinner \texttt{ActivityIndicator} cu
culoarea \texttt{neonPink} (\texttt{\#f472b6}). Abia după finalizarea hidrării
este randat \texttt{RootNavigator}, prevenind orice stare incoherentă de UI în
care componentele ar accesa câmpuri din \texttt{GameState} înainte de
restaurarea lor completă.

% ============================================================
%  CAPITOLUL 5 – GESTIUNEA STĂRII ȘI PERSISTENȚA DATELOR
% ============================================================
\chapter{Gestiunea stării și persistența datelor}

\section{Structura stării globale}

Starea aplicației este definită în \texttt{src/context/GameContext.tsx} ca interfața
TypeScript \texttt{GameState}. Aceasta cuprinde toate datele necesare aplicației (Tabelul~\ref{tab:gamestate}):

\begin{table}[H]
\centering
\caption{Câmpurile principale ale stării globale \texttt{GameState}}
\label{tab:gamestate}
\begin{tabular}{lll}
\toprule
\textbf{Câmp} & \textbf{Tip} & \textbf{Descriere} \\
\midrule
\texttt{profileName} & \texttt{string} & Numele profilului utilizatorului \\
\texttt{coins} & \texttt{number} & Monede virtuale disponibile \\
\texttt{streak} & \texttt{number} & Zile consecutive de concentrare \\
\texttt{bestSessionMinutes} & \texttt{number} & Cea mai lungă sesiune de focus \\
\texttt{totalFocusMinutes} & \texttt{number} & Total minute de focus acumulate \\
\texttt{focusSessions} & \texttt{FocusSession[]} & Ultimele 5 sesiuni de focus \\
\texttt{flashcards} & \texttt{Flashcard[]} & Intrările din jurnal \\
\texttt{chess} & \texttt{ChessProgress} & Statistici șah clasic \\
\texttt{maiaChess} & \texttt{MaiaChessProgress} & Statistici Maia Chess \\
\texttt{sudoku} & \texttt{SudokuProgress} & Statistici și progres Sudoku \\
\texttt{inventory} & \texttt{InventoryItem[]} & Obiectele colectate \\
\texttt{achievements} & \texttt{UserAchievement[]} & Realizări deblocate \\
\texttt{learning} & \texttt{LearningProgress} & Progres module de învățare \\
\texttt{musicEnabled} & \texttt{boolean} & Starea muzicii de fundal \\
\texttt{soundEffectsEnabled} & \texttt{boolean} & Starea efectelor sonore \\
\texttt{totalCoinsEarned} & \texttt{number} & Total monede acumulate de la instalare \\
\texttt{lastSessionDateISO} & \texttt{string?} & Data calendaristică a ultimei sesiuni (ISO) \\
\texttt{arcadeHighScores} & \texttt{\{lanes, reaction: number\}} & Recorduri personale mini-jocuri \\
\texttt{xp} & \texttt{number} & Total puncte de experiență acumulate \\
\texttt{daily} & \texttt{DailyProgress} & Contoare quest-uri zilnice (cu cheie de dată) \\
\texttt{quizStats} & \texttt{QuizStats} & Statistici examene simulate \\
\texttt{puzzleStats} & \texttt{PuzzleStats} & Statistici puzzle-uri de șah rezolvate \\
\bottomrule
\end{tabular}
\end{table}

\section{Sistemul de acțiuni (Reducer)}

Reducer-ul gestionează peste 24 de tipuri de acțiuni, dintre care cele mai importante sunt:

\begin{description}
  \item[\texttt{COMPLETE\_SESSION}] Finalizarea unei sesiuni de focus. Calculează monedele
    câștigate (\textit{baseReward} + \textit{streakBonus}), actualizează streak-ul și
    salvează sesiunea cu datele de geolocation.
  \item[\texttt{ADD\_FLASHCARD / UPDATE\_FLASHCARD / DELETE\_FLASHCARD}] Gestionarea
    intrărilor din jurnal.
  \item[\texttt{COMPLETE\_CHESS\_GAME / COMPLETE\_MAIA\_CHESS\_GAME}] Înregistrarea
    rezultatelor partidelor de șah cu actualizarea statisticilor și acordarea recompenselor.
  \item[\texttt{COMPLETE\_SUDOKU\_GAME}] Finalizarea unui puzzle Sudoku cu acordarea
    monedelor.
  \item[\texttt{UNLOCK\_MAIA\_DIFFICULTY / UNLOCK\_SUDOKU\_DIFFICULTY}] Deblocarea
    nivelelor de dificultate mai ridicate contra cost de monede.
  \item[\texttt{UNLOCK\_ACHIEVEMENT}] Deblocarea unei realizări.
  \item[\texttt{MARK\_EXERCISE\_SOLVED}] Marcarea unui exercițiu ca rezolvat.
  \item[\texttt{COMPLETE\_QUIZ}] Finalizarea unui examen simulat: acordă monede și XP,
    actualizează statisticile și adaugă în progresul de învățare exercițiile rezolvate corect.
  \item[\texttt{COMPLETE\_PUZZLE}] Înregistrarea rezolvării unui puzzle de șah (recompensă
    acordată o singură dată per puzzle).
  \item[\texttt{CLAIM\_QUEST / REFRESH\_DAILY}] Revendicarea unui quest zilnic finalizat,
    respectiv resetarea contoarelor la schimbarea zilei calendaristice.
  \item[\texttt{HYDRATE}] Restaurarea stării din AsyncStorage la pornirea aplicației.
\end{description}

Fiecare acțiune care acordă recompense (sesiuni de focus, victorii la jocuri, exerciții,
examene, puzzle-uri) incrementează simultan câmpul \texttt{xp}, iar acțiunile relevante
actualizează contoarele zilnice utilizate de quest-uri.

\section{Formula de calcul a monedelor din sesiunile de focus}

Monedele câștigate per sesiune de focus sunt calculate conform Ecuației~\ref{eq:coins}:

\begin{equation}
  \text{coinsEarned} = \max(3,\; \lfloor \text{durationMinutes} \times 1{,}5 \rfloor)
  + \left\lfloor \frac{\text{streak}}{3} \right\rfloor
  \label{eq:coins}
\end{equation}

unde \textit{durationMinutes} este durata sesiunii în minute, iar \textit{streak} este
numărul curent de zile consecutive cu cel puțin o sesiune completată.

\section{Sistemul de streak}

Streak-ul este calculat pe baza datei calendaristice (nu a orei):
\begin{itemize}
  \item Dacă ultima sesiune a fost azi: streak rămâne nemodificat;
  \item Dacă ultima sesiune a fost ieri: streak crește cu 1;
  \item Altfel: streak se resetează la 1.
\end{itemize}

\section{Sistemul de experiență (XP) și niveluri}
\label{sec:xp}

Pe lângă economia de monede (resursă consumabilă pentru deblocări), aplicația implementează
un sistem de \textbf{experiență (XP)} cumulativă, ireversibilă, care determină \textbf{nivelul}
jucătorului — un indicator de progres pe termen lung specific jocurilor de rol. Puncte de
experiență sunt acordate pentru aproape fiecare activitate productivă: minute de concentrare,
exerciții rezolvate, victorii la șah/Sudoku, examene simulate, puzzle-uri și revendicarea
quest-urilor zilnice.

XP-ul necesar pentru a avansa de la nivelul $L$ la nivelul $L+1$ crește liniar
($\text{XP\_STEP} \times L$, cu $\text{XP\_STEP} = 100$), astfel încât XP-ul cumulativ
necesar pentru a \emph{atinge} nivelul $L$ este dat de Ecuația~\ref{eq:xp}:

\begin{equation}
  \text{totalXP}(L) = \text{XP\_STEP} \cdot \frac{(L-1)\,L}{2}
  \label{eq:xp}
\end{equation}

Nivelul curent se obține prin inversarea acestei relații (Ecuația~\ref{eq:level}):

\begin{equation}
  L(\text{xp}) = \left\lfloor \frac{1 + \sqrt{1 + 8\,\text{xp}/\text{XP\_STEP}}}{2} \right\rfloor
  \label{eq:level}
\end{equation}

Funcțiile pure de conversie sunt izolate în modulul \texttt{src/constants/leveling.ts}, fiind
reutilizate atât de interfață (bara de progres \texttt{LevelBar}), cât și de sistemul de
realizări. La trecerea într-un nivel superior, aplicația oferă un \textit{feedback} dedicat
(efect sonor de tip fanfară și vibrație), detectat printr-un \texttt{useEffect} care compară
nivelul curent cu cel anterior.

\section{Quest-uri zilnice}
\label{sec:quests}

Sistemul de quest-uri zilnice încurajează revenirea zilnică în aplicație. Trei obiective
sunt disponibile în fiecare zi (Tabelul~\ref{tab:quests}), iar progresul lor este derivat
din contoare zilnice menținute în câmpul \texttt{daily} al stării globale.

\begin{table}[H]
\centering
\caption{Quest-urile zilnice disponibile}
\label{tab:quests}
\begin{tabular}{lcc}
\toprule
\textbf{Obiectiv} & \textbf{Țintă} & \textbf{Recompensă} \\
\midrule
Finalizează o sesiune de focus & 1 & 10 monede + 30 XP \\
Rezolvă 5 exerciții & 5 & 15 monede + 40 XP \\
Câștigă un joc & 1 & 10 monede + 30 XP \\
\bottomrule
\end{tabular}
\end{table}

Structura \texttt{daily} reține cheia zilei curente (\texttt{dateISO}) împreună cu contoarele
de progres (sesiuni de focus, exerciții, victorii) și lista quest-urilor deja revendicate.
La fiecare acțiune relevantă, o funcție utilitară (\texttt{rollDaily}) compară cheia stocată
cu data curentă: dacă ziua s-a schimbat, contoarele și lista de revendicări sunt resetate
automat înainte de incrementare. Astfel, resetarea quest-urilor este \textbf{leneșă}
(\textit{lazy}) — nu necesită un temporizator de fundal, ci se produce la prima interacțiune
din ziua nouă (sau la pornirea aplicației, prin acțiunea \texttt{REFRESH\_DAILY}).

\section{Persistența cu AsyncStorage}

La fiecare modificare a stării, datele sunt serializate JSON și salvate asincron:

\begin{lstlisting}[language=Java, caption={Salvarea și restaurarea stării}]
// Salvare
useEffect(() => {
  if (isHydrated) {
    AsyncStorage.setItem(STORAGE_KEY, JSON.stringify(state));
  }
}, [state, isHydrated]);

// Restaurare la pornire
useEffect(() => {
  AsyncStorage.getItem(STORAGE_KEY).then(raw => {
    if (raw) {
      const parsed = JSON.parse(raw) as GameState;
      dispatch({ type: 'HYDRATE', payload: parsed });
    }
    setIsHydrated(true);
  });
}, []);
\end{lstlisting}

Cheia de stocare este versionată (\texttt{retro-arcade-study-state-v1}) pentru a permite
migrări viitoare de schemă fără a pierde datele existente.

Restaurarea stării (acțiunea \texttt{HYDRATE}) nu este o simplă atribuire a obiectului
deserializat, ci o operație defensivă de îmbinare. Starea salvată este combinată cu
\texttt{initialState} folosit ca bază (\textit{baseline}), astfel încât orice câmp nou,
introdus într-o versiune mai recentă a aplicației și absent din datele salvate anterior, să
primească automat valoarea sa implicită. În plus, reducer-ul revalidează sub-structurile
critice: obiectele de inventar sunt re-rezolvate față de catalogul curent (pentru a prelua
eventualele modificări de denumire sau imagine), iar nivelurile de dificultate deblocate la
șah sunt rederivate din numărul de victorii înregistrate. Acest mecanism protejează
aplicația împotriva datelor parțial corupte sau învechite și asigură că interfața primește
întotdeauna o stare completă și coerentă, indiferent de versiunea în care au fost salvate
inițial datele.

\section{Context API și hook-ul \texttt{useGame}}

Starea și funcțiile de modificare sunt expuse prin \texttt{GameContext} și consumate prin
hook-ul \texttt{useGame()}. Această abordare evită prop drilling și centralizează toată
logica de business.

Alegerea combinației \texttt{Context API + useReducer} în locul unei biblioteci externe de
gestiune a stării (precum Redux sau Zustand) a fost una deliberată: aplicația dispune de o
singură sursă de adevăr (\textit{single source of truth}), iar volumul de logică nu
justifică dependențele și boilerplate-ul suplimentar ale unui store extern. Funcțiile de
acțiune expuse prin context (de exemplu \texttt{completeSession}, \texttt{finishChessMatch},
\texttt{markExerciseSolved}) sunt memoizate, astfel încât referințele lor să rămână stabile
între re-randări și să nu provoace re-randarea inutilă a componentelor consumatoare.
O parte din logica de rezolvare — calculul recompenselor, extragerea aleatorie a obiectelor
de inventar și verificarea pragurilor de deblocare — este încapsulată direct în aceste
funcții, care returnează către interfață obiecte de rezultat (\textit{resolution objects})
folosite ulterior pentru afișarea ferestrelor modale de confirmare. Astfel, componentele de
prezentare rămân subțiri, iar regulile de joc sunt concentrate într-un singur loc, ușor de
testat și de întreținut.

% ============================================================
%  CAPITOLUL 6 – MODULE FUNCȚIONALE
% ============================================================
\chapter{Module funcționale ale aplicației}

\section{Ecranul principal -- Focus Timer (HomeScreen)}

\subsection{Descriere}

Ecranul principal este punctul de intrare în aplicație și cuprinde:
\begin{itemize}
  \item Un mesaj de bun venit cu numele utilizatorului;
  \item Trei badge-uri statistice: monede, focus streak și cea mai lungă sesiune;
  \item Bara de experiență și nivel (componenta \texttt{LevelBar});
  \item Panoul de quest-uri zilnice (componenta \texttt{DailyQuests});
  \item Un player de muzică lo-fi (componenta \texttt{MusicControl});
  \item Timer-ul de concentrare (componenta \texttt{FocusTimer});
  \item Istoricul ultimelor 5 sesiuni de focus cu timestamp și locație.
\end{itemize}

Interfața ecranului principal este prezentată în Figura~\ref{fig:homescreen}, iar
istoricul sesiunilor cu informații de geolocalizare în Figura~\ref{fig:homescreengeo}.

\begin{figure}[H]
\centering
\begin{subfigure}[t]{0.44\textwidth}
  \centering
  \includegraphics[width=\linewidth]{home_screen_with_quests.jpeg}
  \caption{Nivelul curent și misiunile zilnice.\label{fig:homescreen}}
\end{subfigure}
\hfill
\begin{subfigure}[t]{0.44\textwidth}
  \centering
  \includegraphics[width=\linewidth]{home_screen_recent_sessions_with_geolocation.jpeg}
  \caption{Timer-ul de concentrare activ și Istoricul sesiunilor recente cu locația GPS a fiecărei sesiuni.\label{fig:homescreengeo}}
\end{subfigure}
\caption{Ecranul principal al aplicației FocusQuest.}
\end{figure}

\subsection{Componenta FocusTimer}

\texttt{FocusTimer} implementează un cronometru cu numărătoare inversă cu preset-uri
predefinite de 1, 5, 15, 25 și 45 de minute. La finalizarea sesiunii, se apelează
callback-ul \texttt{onSessionComplete(durationMinutes)} care:
\begin{enumerate}
  \item Solicită geolocation-ul curent (dacă permisiunea a fost acordată);
  \item Apelează \texttt{completeSession(minutes, locationData)} din context;
  \item Verifică dacă noi realizări au fost deblocate.
\end{enumerate}

Timer-ul este însoțit de o animație sprite de schelet RPG (\texttt{SkeletonSprite}) care
afișează animații diferite în funcție de stare, controlate printr-o mașină de stare
simplă: pornire/reluare = animație aleatorie, pauză/finalizare = animația \textit{slashing}.

\subsection{Sistemul de animație sprite (\texttt{SkeletonSprite})}

Componenta \texttt{SkeletonSprite} implementează un motor de animație frame-by-frame
cu două optimizări tehnice vizibile în codul sursă:
\begin{enumerate}
  \item \textbf{Buclă de animație bazată pe \texttt{requestAnimationFrame}}
    \textemdash{} în loc de \texttt{setInterval}, componenta folosește
    \texttt{requestAnimationFrame} cu un acumulator de timp delta. Cadrul este
    avansat numai după ce s-a scurs intervalul corect
    ($\lfloor 1000/\text{fps} \rfloor$ ms), asigurând fluiditate independentă
    de încărcarea thread-ului JavaScript;
  \item \textbf{Preîncărcare asincronă a cadrelor} \textemdash{} la schimbarea
    setului de cadre, toate imaginile sunt pre-fetch-uite cu \texttt{Image.prefetch()}
    înainte de pornirea animației, eliminând flickering-ul la prima redare.
\end{enumerate}

Animațiile disponibile sunt descrise în Tabelul~\ref{tab:animations}.

\begin{table}[H]
\centering
\caption{Animațiile sprite ale personajului Skeleton Crusader}
\label{tab:animations}
\begin{tabular}{llll}
\toprule
\textbf{Cheie} & \textbf{Cadre} & \textbf{FPS} & \textbf{Declanșator} \\
\midrule
\texttt{running} & 8 & 12 & Pornire sau reluare a timer-ului \\
\texttt{slashing} & 12 & 15 & Pauză sau finalizarea sesiunii \\
\bottomrule
\end{tabular}
\end{table}

\subsection{Geolocalizarea sesiunilor}

Hook-ul \texttt{useGeolocation} utilizează \texttt{@react-native-community/geolocation}
pentru a obține coordonatele GPS. Acestea sunt trimise API-ului gratuit
\textbf{Nominatim (OpenStreetMap)}~\cite{nominatim} pentru reverse geocoding, obținând numele orașului
și al țării. Aceste date sunt stocate împreună cu sesiunea de focus și afișate în
istoricul sesiunilor.

\subsection{Componenta MusicControl și serviciul MusicPlayer}

Muzica de fundal este gestionată de clasa singleton \texttt{MusicPlayerService} care
utilizează \texttt{react-native-sound}. Sunt definite 5 piese lo-fi:
\begin{itemize}
  \item Cutie Japan Lofi;
  \item Lofi Cafe;
  \item Lofi Girl Ambient;
  \item Lofi Study Calm;
  \item Rainy Lofi City.
\end{itemize}
Utilizatorul poate da play/pause, trece la piesa următoare sau anterioară. Starea
muzicii (activată/dezactivată) este persistată în \texttt{GameState}.

\section{Ecranul Arcade (ArcadeScreen)}

\subsection{Prezentare generală}

Ecranul Arcade grupează cinci mini-jocuri accesibile prin card selector, după cum se
poate observa în Figura~\ref{fig:arcadelist}:
\begin{enumerate}
  \item \textbf{Hyper Lane Defender} — mini-joc reflex;
  \item \textbf{Neon Reaction Pulse} — mini-joc timp de reacție;
  \item \textbf{Cosmic Chess Gauntlet} — șah cu AI bazat pe evaluare de material;
  \item \textbf{Maia AI Chess Challenge} — șah cu modelul Maia ONNX;
  \item \textbf{Sudoku Challenge} — puzzle Sudoku.
\end{enumerate}

\begin{figure}[H]
\centering
\includegraphics[width=0.45\textwidth]{list_of_all_games.jpeg}
\caption{Lista jocurilor disponibile în ecranul Arcade al aplicației FocusQuest.\label{fig:arcadelist}}
\end{figure}

\subsection{Hyper Lane Defender}

Un joc de tip \textit{lane-matching} cu trei culoare (Left, Center, Right) în care
utilizatorul trebuie să apese butonul corespunzător culoarului iluminat. Mecanica
de joc implementată în codul sursă:
\begin{itemize}
  \item \textbf{Viteză adaptivă} \textemdash{} intervalul de schimbare a culoarului
    activ pornește la \textbf{1300 ms} și scade cu \textbf{20 ms} la fiecare
    răspuns corect, până la minimul de \textbf{600 ms}, crescând continuu
    dificultatea pe durata sesiunii;
  \item \textbf{Punctaj} \textemdash{} fiecare apăsare corectă aduce \textbf{8 puncte};
  \item \textbf{Vieți} \textemdash{} jucătorul are \textbf{3 vieți}; o apăsare
    greșită scade o viață, iar la 0 vieți jocul se încheie automat cu
    înregistrarea scorului;
  \item \textbf{Bonus monede} \textemdash{} la înregistrarea scorului se acordă
    $\lfloor \text{scor} / 15 \rfloor$ monede suplimentare, indiferent dacă este
    sau nu un record personal.
\end{itemize}

\subsection{Neon Reaction Pulse}

Un joc de timp de reacție structurat în \textbf{5 runde} cu un maxim de
\textbf{3 greșeli} (apăsări premature) permise. Fluxul unei runde:
\begin{enumerate}
  \item \textbf{Faza de armare (\textit{arming})} \textemdash{} după inițiere,
    un temporizator cu durată aleatorie între \textbf{900 ms și 3100 ms}
    (\texttt{900 + Math.random() * 2200}) precede flash-ul vizual; apăsarea în
    această fereastră este penalizată ca \textit{too soon};
  \item \textbf{Flash și răspuns} \textemdash{} când flash-ul apare, utilizatorul
    trebuie să apese cât mai rapid; latența în milisecunde este înregistrată;
  \item \textbf{Calculul scorului} \textemdash{} scorul per rundă este calculat
    invers proporțional cu latența de reacție.
\end{enumerate}
Jocul se încheie fie după 5 runde de succes, fie la acumularea a 3 greșeli.
Scorul final este convertit în monede ($\lfloor \text{scor} / 15 \rfloor$) care
sunt adăugate automat la economia jucătorului.

\subsection{Jocul de Sudoku (SudokuArena)}

\subsubsection{Implementare}

Componenta \texttt{SudokuArena} implementează un joc complet de Sudoku cu:
\begin{itemize}
  \item Trei niveluri de dificultate: \textit{easy}, \textit{medium}, \textit{expert};
  \item Puzzle-uri preîncărcate din constantele \texttt{EASY\_PUZZLES},
    \texttt{MEDIUM\_PUZZLES}, \texttt{EXPERT\_PUZZLES};
  \item Un algoritm de backtracking pentru validarea și rezolvarea puzzle-urilor.
\end{itemize}

\subsubsection{Algoritmul de backtracking}

Puzzle-urile Sudoku sunt rezolvate prin backtracking recursiv. Algoritmul:
\begin{enumerate}
  \item Caută prima celulă goală (\textit{findEmpty});
  \item Încearcă numerele de la 1 la 9;
  \item Verifică validitatea plasării față de rând, coloană și sub-carou $3\times3$;
  \item Dacă nu există număr valid, revine la celula anterioară.
\end{enumerate}

\begin{lstlisting}[language=Java, caption={Algoritmul de backtracking Sudoku}]
const solveSudoku = (grid: (number | null)[][]): boolean => {
  const findEmpty = (): [number, number] | null => {
    for (let row = 0; row < 9; row++)
      for (let col = 0; col < 9; col++)
        if (grid[row][col] === null) return [row, col];
    return null;
  };

  const isValid = (num: number, row: number, col: number): boolean => {
    for (let x = 0; x < 9; x++)
      if (grid[row][x] === num || grid[x][col] === num) return false;
    const sr = Math.floor(row / 3) * 3;
    const sc = Math.floor(col / 3) * 3;
    for (let i = 0; i < 3; i++)
      for (let j = 0; j < 3; j++)
        if (grid[sr + i][sc + j] === num) return false;
    return true;
  };

  const empty = findEmpty();
  if (!empty) return true; // Rezolvat!
  const [r, c] = empty;
  for (let num = 1; num <= 9; num++) {
    if (isValid(num, r, c)) {
      grid[r][c] = num;
      if (solveSudoku(grid)) return true;
      grid[r][c] = null;
    }
  }
  return false;
};
\end{lstlisting}

\subsubsection{Sistemul de penalizare}

Aspectul vizual al unui joc de Sudoku în desfășurare este prezentat în
Figura~\ref{fig:sudoku}.

\begin{figure}[H]
\centering
\includegraphics[width=0.45\textwidth]{sudoku_game.jpeg}
\caption{Jocul de Sudoku în aplicația FocusQuest, cu grila de 9×9 și bara de progres.\label{fig:sudoku}}
\end{figure}

Utilizatorul are maxim 3 greșeli per partidă. La depășirea limitei, jocul se încheie.
Celulele introduse greșit sunt marcate vizual, iar la corecție marcajul dispare.

\subsubsection{Recompense Sudoku}

Recompensele și costurile de deblocare pe dificultăți sunt prezentate în Tabelul~\ref{tab:sudoku_rewards}.

\begin{table}[H]
\centering
\caption{Recompense Sudoku pe dificultăți}
\label{tab:sudoku_rewards}
\begin{tabular}{lcc}
\toprule
\textbf{Dificultate} & \textbf{Cost deblocare} & \textbf{Recompensă finalizare} \\
\midrule
Easy & Gratuit & 10 monede \\
Medium & Gratuit & 20 monede \\
Expert & 7 monede & 40 monede \\
\bottomrule
\end{tabular}
\end{table}

\section{Modulul de Jurnal (FlashcardsScreen)}

\subsection{Descriere}

Modulul de jurnal permite utilizatorului să creeze intrări de tip \textit{flashcard} care
conțin:
\begin{itemize}
  \item Titlu și descriere text;
  \item O imagine opțională (din galerie sau cameră foto);
  \item O stare emoțională (\textit{mood}): reflectiv, recunoscător, energizat, curios,
    victorios.
\end{itemize}

Ecranul jurnalului cu mai multe intrări existente este ilustrat în Figura~\ref{fig:journal}.

\begin{figure}[H]
\centering

\includegraphics[width=0.45\textwidth]{journal_screen_with_entries.jpeg}
\caption{Ecranul jurnalului digital cu lista intrărilor utilizatorului și stările emoționale asociate.\label{fig:journal}}
\end{figure}

\subsection{Structura datelor Flashcard}

\begin{lstlisting}[language=Java, caption={Structura tipului Flashcard}]
export type Flashcard = {
  id: string;
  title: string;
  description: string;
  imageBase64?: string;
  createdAt: number;
  lastReviewedAt?: number;
  favorite?: boolean;
  mood?: JournalMood;
};
\end{lstlisting}

Imaginile sunt stocate ca șiruri Base64 direct în starea aplicației (persistent în
AsyncStorage), eliminând necesitatea unui sistem de fișiere separat.

\subsection{Clasificarea AI a imaginilor}

La adăugarea unei imagini, hook-ul \texttt{useOnnxClassifier} este invocat automat pentru
a rula inferența ResNet-50 și a sugera un titlu pentru intrarea de jurnal. Utilizatorul
poate accepta sugestia sau poate introduce manual titlul. În Figura~\ref{fig:resnetjournal}
este prezentată clasificarea automată a unei imagini din jurnal de către modelul ResNet-50.

\begin{figure}[H]
\centering
\includegraphics[width=0.45\textwidth]{journal_resnet_recognizing_image.jpeg}
\caption{Clasificarea automată a imaginii din jurnal prin modelul ResNet-50 INT8 rulat local.\label{fig:resnetjournal}}
\end{figure}

\section{Ecranul de Profil (ProfileScreen)}

Ecranul de profil permite:
\begin{itemize}
  \item Editarea numelui de profil (câmpul \textit{Call sign});
  \item Vizualizarea statisticilor generale: monede, streak, minute de focus;
  \item Vizualizarea top 6 sesiuni de focus sortate descendent după durată;
  \item Navigarea spre Inventar, Realizări și Statistici.
\end{itemize}

\section{Ecranul de Statistici (StatisticsScreen)}

Dashboard-ul de statistici agregă indicatorii de performanță din toate modulele aplicației
(Figura~\ref{fig:statistics}):
\begin{enumerate}
  \item \textbf{Nivel} \textemdash{} bara de XP și nivelul curent (componenta \texttt{LevelBar});
  \item \textbf{Focus Sessions} \textemdash{} total sesiuni, sesiuni \^{\i}n ultimele 7/30 zile,
    durat\u{a} medie, durat\u{a} maxim\u{a}, total minute de focus;
  \item \textbf{Chess Statistics} — total jocuri, victorii, rată de succes, detalii pe dificultăți;
  \item \textbf{Maia AI Chess} — jocuri, victorii și rată de succes pe cele trei niveluri Elo;
  \item \textbf{Sudoku} — puzzle-uri jucate și finalizate, pe dificultăți;
  \item \textbf{Învățare} — exerciții rezolvate pe discipline, acuratețea la examenele
    simulate și numărul de puzzle-uri de șah rezolvate;
  \item \textbf{Journal Insights} — total intrări, favorite, starea emoțională predominantă;
  \item \textbf{Economy} — total monede câștigate;
  \item \textbf{Productivity Insights} — ora cu cele mai multe sesiuni de focus completate.
\end{enumerate}

Față de versiunea inițială (care raporta doar șahul clasic), dashboard-ul a fost extins pentru
a acoperi integral cele trei jocuri inteligente, progresul educațional și sistemul de niveluri.

\begin{figure}[H]
\centering
\includegraphics[width=0.45\textwidth]{statistics_Screen.jpeg}
\caption{Ecranul de statistici cu indicatorii de performanță ai utilizatorului.\label{fig:statistics}}
\end{figure}

\section{Ecranul de Realizări (AchievementsScreen)}

Sistemul de realizări cuprinde \textbf{42 de realizări} organizate în 12 categorii:
focus, șah clasic, Maia Chess, Sudoku, jurnal, monede, arcade, inventar, învățare, examen
simulat (\textit{quiz}), puzzle-uri și niveluri. Ultimele trei categorii au fost adăugate
odată cu noile module de gamificare și educație. Ecranul de realizări este prezentat în
Figura~\ref{fig:achievements}. Exemple din fiecare categorie sunt listate în
Tabelul~\ref{tab:achievements}.

\begin{figure}[H]
\centering
\includegraphics[width=0.45\textwidth]{achievements_screen.jpeg}
\caption{Ecranul de realizări cu insigne deblocate și progress per categorie.\label{fig:achievements}}
\end{figure}

\begin{table}[H]
\centering
\caption{Exemple de realizări pe categorii}
\label{tab:achievements}
\begin{tabular}{lll}
\toprule
\textbf{ID} & \textbf{Titlu} & \textbf{Condiție} \\
\midrule
first-session & First Focus & Prima sesiune completată \\
marathon-runner & Marathon Runner & Sesiune de $\geq 90$ min \\
night-owl & Night Owl & 10 sesiuni după ora 22 \\
chess-legend & Chess Legend & 50 victorii la șah \\
maia-champion & AI Champion & Victorie vs Maia pe Hard \\
sudoku-expert & Sudoku Expert & 10 victorii Sudoku totale \\
book-master & Book Master & Toate exercițiile rezolvate \\
streak-keeper & Streak Keeper & Streak de $\geq 7$ zile \\
completionist & Completionist & 78 obiecte (48 loot + 30 săbii) \\
journaler & Journaler & 30 de intrări în jurnal \\
quiz-ace & Top of the Class & Scor de 100\% la un examen simulat \\
checkmate-virtuoso & Checkmate Virtuoso & 30 de puzzle-uri de șah rezolvate \\
seasoned-adventurer & Seasoned Adventurer & Atingerea nivelului 10 \\
\bottomrule
\end{tabular}
\end{table}

\subsection{Mecanismul de deblocare automată}

Spre deosebire de sistemele care necesită o verificare declanșată manual, deblocarea
realizărilor în FocusQuest este complet automată. Un efect lateral (\texttt{useEffect})
prevăzut cu un mecanism de \textit{debounce} de 500\,ms monitorizează modificările stării
globale și reevaluează, la fiecare schimbare relevantă, condițiile tuturor celor 42 de
realizări. Întârzierea de 500\,ms previne reevaluările redundante atunci când mai multe
câmpuri ale stării se modifică în rafală — de exemplu, la finalizarea unei sesiuni de
concentrare, care actualizează simultan monedele, streak-ul, totalul de minute și istoricul
sesiunilor. În momentul în care numărul de realizări deblocate crește, aplicația oferă un
feedback haptic prin vibrație (\texttt{Vibration.vibrate}), confirmând tactil progresul
utilizatorului fără a întrerupe fluxul curent al interfeței. Această abordare reactivă
garantează că o realizare este acordată imediat ce condiția sa este îndeplinită, indiferent
de ecranul pe care se află utilizatorul în acel moment.

\section{Ecranul de Inventar (InventoryScreen)}

Inventarul stochează obiectele câștigate din jocuri. Catalogul conține:
\begin{itemize}
  \item \textbf{48 de obiecte de tip \textit{loot}}: cristale, pergamente, elixiruri,
    relicve;
  \item \textbf{30 de arme de tip \textit{sword}}: săbii, suliți, topoare cu tematică
    fantasy și sci-fi.
\end{itemize}

Obiectele sunt acordate aleatoriu la finalizarea partidelor de șah pe dificultăți medii
sau ridicate. Inventarul este afișat într-o grilă cu 3 coloane, sortată după data
obținerii, după cum se observă în Figura~\ref{fig:inventory}.

\begin{figure}[H]
\centering
\includegraphics[width=0.45\textwidth]{inventory_screen.jpeg}
\caption{Ecranul de inventar cu obiectele colectate afișate în grilă.\label{fig:inventory}}
\end{figure}

\section{Modulul de Învățare (LearningHubScreen)}

\subsection{Structura conținutului}

Modulul de învățare oferă un total de \textbf{673 de exerciții matematice} structurate pe
trei discipline (Tabelul~\ref{tab:learning}): 338 de algebră, 246 de analiză și 89 de
trigonometrie. Numărul de exerciții pe disciplină și totalul general sunt calculate
\textbf{dinamic} la rulare din fișierele de conținut, astfel încât adăugarea de noi exerciții
nu necesită modificarea codului.

\begin{table}[H]
\centering
\caption{Disciplinele modulului de învățare}
\label{tab:learning}
\begin{tabular}{llll}
\toprule
\textbf{Disciplină} & \textbf{Subteme} & \textbf{Cost deblocare} & \textbf{Prefix ID} \\
\midrule
Algebră & Inecuații, valori absolute, ecuații, teoria nr. & 100 monede & AL \\
Analiză & Limite, derivate, continuitate, asimptote & 100 monede & AM \\
Trigonometrie & Unghiuri, identități, ecuații, funcții & 100 monede & TG \\
\bottomrule
\end{tabular}
\end{table}

Interfața hub-ului de învățare cu cele trei discipline disponibile este prezentată în
Figura~\ref{fig:learninghub}.

\begin{figure}[H]
\centering

\includegraphics[width=0.45\textwidth]{learning_hub_with_quiz.jpeg}
\caption{Hub-ul de învățare cu cele trei discipline matematice și starea de deblocare a fiecăreia.\label{fig:learninghub}}
\end{figure}

Datele sunt stocate în fișiere JSON (\texttt{assets/learning/algebra.json},
\texttt{analysis.json}, \texttt{trigonometry.json}) cu structura:

\begin{lstlisting}[language=Java, caption={Structura unui exercițiu de învățare}]
export type LearningExercise = {
  id: string;         // ex: "AL001"
  problem: string;    // formularea problemei (poate contine LaTeX)
  answers: LearningAnswer[]; // 4-5 variante de raspuns
  correctAnswer: string;     // ID-ul raspunsului corect
  solution: string;   // solutia pas-cu-pas (LaTeX)
};
\end{lstlisting}

\subsection{Randarea formulelor matematice (MathRenderer)}

Componenta \texttt{MathRenderer} este construită pe un \texttt{WebView} care injectează
biblioteca KaTeX pentru randarea formulelor LaTeX. Caracteristici tehnice:
\begin{itemize}
  \item Suportă atât formule inline (\texttt{\$...\$}) cât și formule bloc (\texttt{\$\$...\$\$});
  \item Definește macrouri pentru notațiile românești: \texttt{\textbackslash tg},
    \texttt{\textbackslash ctg}, \texttt{\textbackslash arctg}, \texttt{\textbackslash arcctg};
  \item Înălțimea WebView-ului se ajustează dinamic prin mesaje
    \texttt{postMessage} de la KaTeX;
  \item Biblioteca KaTeX este încărcată din CDN-ul public \texttt{jsdelivr}, ceea ce face ca
    randarea formulelor să fie singura componentă a aplicației dependentă de conectivitate.
\end{itemize}

\subsubsection{Singura dependență de conectivitate}
\label{sec:online-dependency}

Spre deosebire de restul aplicației, componenta \texttt{MathRenderer} încarcă biblioteca
KaTeX în mod dinamic dintr-o rețea de distribuție de conținut (CDN), motiv pentru care
randarea formulelor matematice din modulul de învățare reprezintă \textbf{singura
funcționalitate care necesită o conexiune activă la internet}. Această abordare a fost
preferată în detrimentul împachetării locale a bibliotecii (\textit{bundling}) pentru a
păstra dimensiunea pachetului aplicației redusă și pentru a beneficia de versiunea
optimizată și cache-uită a CDN-ului. O direcție imediată de îmbunătățire (Capitolul~9)
este împachetarea fișierelor KaTeX direct în resursele aplicației, ceea ce ar elimina
și această ultimă dependență, transformând FocusQuest într-o aplicație complet
funcțională \textit{offline}. Restul funcționalităților — inferența ONNX (ResNet-50 și
Maia Chess), persistența datelor, jocurile, timer-ul de concentrare și sistemul de
gamificare — operează în întregime local, fără niciun apel de rețea. Singura excepție
parțială o constituie geolocalizarea: obținerea coordonatelor GPS este locală, însă
transformarea acestora în denumirea localității (\textit{reverse geocoding}) folosește
serviciul Nominatim; în absența conexiunii, această etapă eșuează elegant, sesiunea
fiind salvată cu coordonatele numerice brute, fără a bloca funcționalitatea.

\subsection{Fluxul de rezolvare a exercițiilor}

\begin{enumerate}
  \item Utilizatorul selectează o disciplină și un exercițiu;
  \item Problema și cele 4–5 variante de răspuns sunt afișate cu MathRenderer;
  \item Utilizatorul selectează un răspuns — imediat se afișează feedback vizual
    (verde = corect, roșu = greșit);
  \item Dacă răspunsul este corect, exercițiul este marcat ca rezolvat în
    \texttt{learning.solvedExercises};
  \item Utilizatorul poate afișa soluția pas-cu-pas independent de răspuns.
\end{enumerate}

În Figura~\ref{fig:exercise} și Figura~\ref{fig:exercisesol} sunt prezentate
ecranul unui exercițiu în așteptarea răspunsului, respectiv același exercițiu
cu soluția pas-cu-pas afișată.

\begin{figure}[H]
\centering
\begin{subfigure}[b]{0.44\textwidth}
  \centering
  \includegraphics[width=\textwidth]{learning_Screen_exercise.jpeg}
  \caption{Exercițiu cu variantele de răspuns.\label{fig:exercise}}
\end{subfigure}
\hfill
\begin{subfigure}[b]{0.44\textwidth}
  \centering
  \includegraphics[width=\textwidth]{learning_screen_exercise_with_solution.jpeg}
  \caption{Același exercițiu cu soluția desfășurată.\label{fig:exercisesol}}
\end{subfigure}
\caption{Fluxul de rezolvare a unui exercițiu matematic în modulul de învățare FocusQuest.}
\end{figure}

\section{Examenul simulat (QuizScreen)}

Pe lângă practica liberă de exerciții, modulul de învățare include un mod de \textbf{examen
simulat} (\textit{mock exam}), inspirat de structura probelor de admitere de tip grilă.
La fiecare examen, aplicația extrage aleatoriu \textbf{12 întrebări} — câte \textbf{4 din
fiecare disciplină} (algebră, analiză, trigonometrie) — și le amestecă, astfel încât
disciplinele să se intercaleze.

% TODO: captură de ecran QuizScreen (întrebare + variante de răspuns)
% \begin{figure}[H]\centering\includegraphics[width=0.45\textwidth]{quiz_screen.jpeg}
% \caption{Ecranul de examen simulat.\label{fig:quiz}}\end{figure}

Fluxul unui examen:
\begin{enumerate}
  \item Utilizatorul selectează un răspuns; se afișează imediat \textit{feedback} vizual și
    sonor (corect/greșit) și se trece la întrebarea următoare;
  \item La final, un ecran de rezultate prezintă scorul (de exemplu $9/12$), procentajul,
    monedele și XP-ul câștigate;
  \item Răspunsurile greșite pot fi revizuite, fiecare împreună cu varianta corectă, randate
    cu \texttt{MathRenderer}.
\end{enumerate}

Recompensele sunt proporționale cu performanța: \textbf{5 monede} și \textbf{10 XP} pentru
fiecare răspuns corect, plus un bonus de \textbf{20 de monede} pentru un scor perfect.
Întrebările răspunse corect sunt totodată marcate ca rezolvate în progresul de învățare,
contribuind la realizările de tip \textit{scholar}. Refolosirea componentei \texttt{MathRenderer}
asigură afișarea consecventă a formulelor matematice atât în practica liberă, cât și în examen.

\section{Sistemul de efecte sonore}
\label{sec:sfx}

Aplicația oferă \textbf{efecte sonore} scurte de tip \textit{chiptune} (8-bit), potrivite
esteticii retro-arcade, pentru momentele-cheie de interacțiune: monedă câștigată, răspuns
corect, răspuns greșit, victorie și avansare în nivel. Spre deosebire de muzica de fundal
(fișiere lo-fi externe), efectele sonore au fost \textbf{sintetizate programatic} ca fișiere
WAV (unde pătrate și sinusoidale cu anvelope scurte), ceea ce elimină orice problemă de
licențiere și menține dimensiunea totală sub 80\,KB.

Redarea este gestionată de un serviciu de tip \textit{singleton} (\texttt{src/services/SoundEffects.ts})
care preîncarcă efectele prin \texttt{react-native-sound} și respectă preferința utilizatorului
(comutatorul \texttt{soundEffectsEnabled} din starea globală). Efectele sunt declanșate din
locurile relevante: răspunsurile la exerciții și examene, revendicarea quest-urilor, deblocarea
realizărilor și avansarea în nivel.

% ============================================================
%  CAPITOLUL 7 – INTEGRAREA MODELELOR DE INTELIGENȚĂ ARTIFICIALĂ
% ============================================================
\chapter{Integrarea modelelor de Inteligență Artificială}

\section{Arhitectura generală de inferență ONNX}

Ambele modele AI utilizate în aplicație sunt în format \textbf{ONNX (Open Neural Network
Exchange)} și sunt rulate local prin pachetul \texttt{onnxruntime-react-native}.
Acest lucru asigură:
\begin{itemize}
  \item Funcționalitate completă offline;
  \item Latență redusă (fără round-trip de rețea);
  \item Confidențialitatea datelor utilizatorului (imaginile nu sunt trimise pe servere
    externe).
\end{itemize}

Modelele sunt incluse ca resurse statice ale aplicației:
\begin{itemize}
  \item Pe \textbf{Android}: copierea din directorul \texttt{assets/models/} în
    \texttt{CachesDirectoryPath} la prima rulare (prin \texttt{react-native-fs});
  \item Pe \textbf{iOS}: accesate direct din \texttt{MainBundlePath}.
\end{itemize}

\section{Modelul ResNet-50 pentru clasificarea imaginilor}

\subsection{Descriere model}

\textbf{ResNet-50}~\cite{resnet} (Residual Network cu 50 de straturi) este o arhitectură CNN
(Convolutional Neural Network) dezvoltată de Microsoft Research, câștigătoare a
competiției ImageNet~\cite{imageNet} Large Scale Visual Recognition Challenge (ILSVRC) 2015.

Versiunea utilizată este \textbf{ResNet-50-v1-7 quantizat INT8} din ONNX Model Zoo~\cite{onnxmodelzoo},
cu dimensiune redusă față de versiunea float32 originală. Modelul clasifică imaginile
în \textbf{1000 de clase ImageNet}.

Parametrii modelului:
\begin{itemize}
  \item Dimensiune intrare: $1 \times 3 \times 224 \times 224$ (batch, canale RGB, înălțime, lățime);
  \item Dimensiune ieșire: $1 \times 1000$ (logit-uri pentru 1000 de clase);
  \item Normalizare intrare: medie $\mu = [0.485, 0.456, 0.406]$, deviație standard
    $\sigma = [0.229, 0.224, 0.225]$ (standardul ImageNet).
\end{itemize}

\subsection{Hook-ul \texttt{useOnnxClassifier}}

Hook-ul \texttt{useOnnxClassifier} gestionează ciclul de viață al sesiunii ONNX:
\begin{enumerate}
  \item \textbf{Încărcare model} — asincron la montarea componentei;
  \item \textbf{Preprocesare imagine} — decodificare JPEG, redimensionare la $224\times224$,
    normalizare și construirea tensorului;
  \item \textbf{Inferență} — \texttt{session.run(feeds)};
  \item \textbf{Postprocesare} — aplicarea \textit{softmax} și selectarea top-K clase.
\end{enumerate}

\subsection{Preprocesarea imaginilor}

Preprocesarea implementată reproduce exact pipeline-ul standard ImageNet:

\begin{lstlisting}[language=Java, caption={Preprocesarea imaginii pentru ResNet-50}]
function preprocessImage(base64: string): Tensor {
  const buffer = Buffer.from(base64, 'base64');
  const decoded = decodeJpeg(buffer, { useTArray: true });
  const resized = resizeToSquare(decoded, 224);
  const tensorData = new Float32Array(3 * 224 * 224);

  // Normalizare ImageNet
  const mean = [0.485, 0.456, 0.406];
  const std  = [0.229, 0.224, 0.225];

  for (let y = 0; y < 224; y++) {
    for (let x = 0; x < 224; x++) {
      const idx = y * 224 + x;
      const si  = idx * 4; // sursa RGBA
      // Format planar: R, G, B in blocuri separate
      tensorData[idx]               = (resized.data[si]   / 255 - mean[0]) / std[0];
      tensorData[idx + 224 * 224]   = (resized.data[si+1] / 255 - mean[1]) / std[1];
      tensorData[idx + 224 * 224*2] = (resized.data[si+2] / 255 - mean[2]) / std[2];
    }
  }
  return new Tensor('float32', tensorData, [1, 3, 224, 224]);
}
\end{lstlisting}

\subsection{Funcția softmax și selectarea predicțiilor}

Logit-urile de ieșire ale modelului sunt transformate în probabilități prin funcția
softmax numerică stabilă (Ecuația~\ref{eq:softmax}):
\begin{equation}
  \text{softmax}(z_i) = \frac{e^{z_i - \max(\mathbf{z})}}{\sum_{j=1}^{K} e^{z_j - \max(\mathbf{z})}}
  \label{eq:softmax}
\end{equation}

Predicțiile sunt sortate descendent după probabilitate și returnate primele K (implicit $K=3$).
Etichetele textuale corespunzătoare indicilor de clasă sunt extrase din fișierul
\texttt{classes.json}, care conține denumirile celor 1000 de clase ImageNet.

\subsection{Integrarea în fluxul jurnalului}

În contextul modulului de jurnal, rezultatul clasificării nu este impus utilizatorului, ci
oferit ca \textbf{sugestie}. La selectarea unei imagini din galerie sau de la cameră,
imaginea este convertită în Base64, decodificată, preprocesată și transmisă modelului, iar
primele trei etichete cu probabilitatea cea mai mare sunt propuse ca posibil titlu al
intrării de jurnal. Utilizatorul poate accepta una dintre sugestii printr-o singură atingere
sau poate ignora complet propunerile și introduce manual un titlu. Această proiectare
păstrează controlul deplin al utilizatorului asupra conținutului, folosind modelul AI ca
instrument de asistență, nu ca mecanism de etichetare automată obligatorie. În eventualitatea
unei erori de inferență sau a unei imagini ce nu poate fi decodificată, fluxul de creare a
intrării continuă neîntrerupt, fără sugestii, evitând blocarea interfeței.

\section{Modelul Maia Chess pentru jocul de șah cu AI}

\subsection{Descriere model Maia Chess}

\textbf{Maia Chess}~\cite{maiacChess} este un proiect de cercetare open-source al Universității Cornell
(Reid McIlroy-Young et al., 2020) care antrenează rețele neuronale să joace șah imitând
stilul uman la diverse niveluri Elo. Codul sursă și modelele sunt disponibile pe
GitHub~\cite{maiaGitHub}. Spre deosebire de motoarele tradiționale (Stockfish,
Leela Chess Zero) care maximizează forța de joc, Maia optimizează \textit{acordul cu
mutările umane} la un Elo dat.

Modelul utilizat este \textbf{maia\_rapid.onnx} — un model unic (rapid model) care
acoperă intervalul Elo 1100–1900. Nivelul de joc este controlat prin doi parametri
de intrare suplimentari: \texttt{elo\_self} și \texttt{elo\_oppo}.

\subsection{Arhitectura modelului}

Modelul Maia are \textbf{3 intrări}:
\begin{enumerate}
  \item \texttt{boards} — tensor de tip \texttt{float32} cu forma $[1, 18, 8, 8]$
    reprezentând starea tablei;
  \item \texttt{elo\_self} — tensor de tip \texttt{int64} cu forma $[1]$, categoria Elo
    a modelului (0–10);
  \item \texttt{elo\_oppo} — tensor de tip \texttt{int64} cu forma $[1]$, categoria Elo
    a adversarului.
\end{enumerate}

Și o \textbf{ieșire}:
\begin{itemize}
  \item \texttt{logits\_maia} — tensor de tip \texttt{float32} cu forma $[1, 1880]$,
    reprezentând scorul pentru fiecare mutare din dicționarul oficial de 1880 mutări
    UCI al proiectului Maia.
\end{itemize}

\subsection{Codificarea stării tablei (18 canale)}

Starea tablei este codificată ca un tensor 3D cu 18 canale binare,
descrise în Tabelul~\ref{tab:maia_channels}:

\begin{table}[H]
\centering
\caption{Structura celor 18 canale de intrare ale modelului Maia}
\label{tab:maia_channels}
\begin{tabular}{cl}
\toprule
\textbf{Canal(e)} & \textbf{Conținut} \\
\midrule
0–5 & Piese albe: P(pion), N(cal), B(nebun), R(tură), Q(regină), K(rege) \\
6–11 & Piese negre: p, n, b, r, q, k \\
12 & Rândul la mutare (1=alb, 0=negru) \\
13 & Rocada regală albă (K) \\
14 & Rocada mare albă (Q) \\
15 & Rocada regală neagră (k) \\
16 & Rocada mare neagră (q) \\
17 & Câmpul en passant \\
\bottomrule
\end{tabular}
\end{table}

\subsection{Tehnica de oglindire FEN}

O problemă critică în implementare este că modelul Maia a fost antrenat întotdeauna
din perspectiva albului. Prin urmare, când negrul este la mutare, tabla trebuie
\textbf{oglindită} înainte de a fi transmisă modelului:

\begin{enumerate}
  \item Rândurile FEN sunt inversate (rândul 8 devine rândul 1 și invers);
  \item Culorile pieselor sunt interschimbate (litere mari $\leftrightarrow$ litere mici);
  \item Rândul la mutare este interschimbat (w $\leftrightarrow$ b);
  \item Drepturile de rocadă sunt oglindite (K $\leftrightarrow$ k, Q $\leftrightarrow$ q);
  \item Câmpul en passant este oglindit (rândul $r \rightarrow 9-r$).
\end{enumerate}

Dacă board-ul a fost oglindit, mutarea selectată de model trebuie și ea oglindită
înapoi la coordonatele originale înainte de a fi executată.

\subsection{Selectarea mutării optime}

Algoritmul de selecție:
\begin{enumerate}
  \item Se generează lista tuturor mutărilor legale cu \texttt{chess.js};
  \item Fiecare mutare legală este convertită în format UCI (cu oglindire dacă e
    necesar) și căutată în \texttt{all\_moves.json} (1880 intrări);
  \item Logit-ul corespunzător este extras din ieșirea modelului;
  \item Se aplică softmax \textbf{doar} pe mutările legale (nu pe toți cei 1880
    candidați) pentru a preveni dispersia probabilistică;
  \item Se returnează mutarea cu probabilitatea maximă.
\end{enumerate}

\begin{lstlisting}[language=Java, caption={Selecția mutării optime Maia}]
const logits = moveScores.map(ms => ms.logit);
const probs  = softmax(logits); // Softmax DOAR pe mutarile legale!

const movesWithProbs = moveScores.map((ms, i) => ({
  ...ms, prob: probs[i]
}));
movesWithProbs.sort((a, b) => b.prob - a.prob);
return movesWithProbs[0].move; // Mutarea cu probabilitate maxima
\end{lstlisting}

\subsection{Niveluri de dificultate Maia}

Cele trei niveluri de dificultate disponibile sunt prezentate în Tabelul~\ref{tab:maia_diff}.

\begin{table}[H]
\centering
\caption{Niveluri de dificultate Maia Chess}
\label{tab:maia_diff}
\begin{tabular}{llcc}
\toprule
\textbf{Dificultate} & \textbf{Elo AI} & \textbf{Cost deblocare} & \textbf{Recompensă victorie} \\
\midrule
Apprentice (Easy) & $\sim$1100 & Gratuit & 15 monede \\
Adept (Medium) & $\sim$1300 & 50 monede & 30 monede \\
Master (Hard) & $\sim$1500 & 100 monede & 50 monede \\
\bottomrule
\end{tabular}
\end{table}

\textbf{Perspectiva jucătorului:} implicit, jucătorul controlează piesele \textbf{albe},
iar Maia joacă piesele \textbf{negre}. ELO-ul jucătorului este setat la valoarea fixă de
\textbf{1600} (\texttt{elo\_oppo = 1600}) în codul sursă, indiferent de nivelul real al
utilizatorului. Acest parametru influențează distribuția de probabilitate a mutărilor
modelului Maia: la un \texttt{elo\_oppo} ridicat, modelul va tinde să aleagă mutări mai
solide, adaptându-și stilul ca și cum ar juca împotriva unui jucător de nivel
intermediar-avansat.

\subsection{Jocul cu piesele negre}
\label{sec:play-black}

Atât în modul Maia, cât și în șahul clasic, jucătorul poate alege să joace cu piesele
\textbf{negre} printr-un buton de inversare (\textit{flip}). În acest caz tabla este rotită
(prin proprietatea \texttt{orientation} a componentei \texttt{ChessBoard}), iar AI-ul preia
piesele albe și mută primul. Modelul Maia nu necesită antrenare separată pentru a juca cu
albul: tehnica de oglindire FEN descrisă anterior aduce întotdeauna poziția în perspectiva
albului înainte de inferență, deci motorul funcționează identic indiferent de culoarea
jucată. Pentru șahul clasic, funcția de evaluare a fost generalizată astfel încât scorul să
fie calculat din perspectiva culorii controlate de AI (maximizare pentru alb, respectiv
pentru negru).

Interfața vizuală a partidei de șah contra modelului Maia este ilustrată în
Figura~\ref{fig:chess}.

\begin{figure}[H]
\centering
\includegraphics[width=0.45\textwidth]{chess_with_ai.jpeg}
\caption{Partida de șah contra modelului Maia Chess cu tabla de joc în stil retro-neon.\label{fig:chess}}
\end{figure}

\subsection{Cartea de deschideri}

Primele mutări ale AI-ului sunt extrase dintr-o carte de deschideri fixă definită
în \texttt{src/constants/chessOpenings.ts}. La începutul fiecărei partide, aplicația
alege aleatoriu una dintre cele 10 deschideri predefinite, pe care modelul Maia
o joacă pentru primele \textbf{două mutări ale negrului}; ulterior, rețeaua neurală
preia controlul. Deschiderile disponibile sunt:

\begin{itemize}
  \item Apărarea Siciliană (\textit{Sicilian Defense});
  \item Apărarea Franceză (\textit{French Defense});
  \item Apărarea Caro-Kann;
  \item Apărarea Pirc;
  \item Apărarea Modernă;
  \item Apărarea Alekhine;
  \item Apărarea Scandinavă;
  \item Apărarea Nimzowitsch;
  \item Deschidere de pion-rege (\textit{King's Pawn});
  \item Configurația Indiană de Rege (\textit{King's Indian Setup}).
\end{itemize}

Motivația cărții de deschideri este dublă: pe de o parte introduce \textbf{varietate}
(altfel modelul ar repeta aproape aceleași mutări în deschidere), iar pe de altă parte
evită predicțiile instabile ale rețelei în primele mutări, când poziția este aproape
simetrică.

Atunci când jucătorul alege piesele negre (Secțiunea~\ref{sec:play-black}), Maia joacă cu
albul și folosește o \textbf{a doua carte de deschideri}, de data aceasta cu deschideri de
alb (de exemplu \textit{King's Knight}, \textit{Queen's Gambit}, \textit{English Opening},
\textit{Réti}, \textit{London System}). A doua mutare din fiecare deschidere de alb a fost
aleasă astfel încât să rămână legală indiferent de răspunsul negrului — proprietate
verificată automat cu \texttt{chess.js} pentru toate combinațiile, pentru a garanta că AI-ul
nu recurge niciodată la o mutare aleatoare de rezervă în deschidere.

\section{Jocul de Șah Clasic (NeonChessArena)}

Pe lângă Maia Chess, aplicația include un joc de șah clasic denumit \textbf{Cosmic Chess
Gauntlet} (componenta \texttt{NeonChessArena}). Motorul AI este implementat cu o funcție
de evaluare bazată pe valoarea materialului fiecărei piese, utilizând algoritmul
minimax~\cite{minimax}:

\begin{itemize}
  \item \textbf{Easy} — mutare aleatoare din lista mutărilor legale;
  \item \textbf{Normal} — evaluare la adâncimea 1: se alege mutarea care minimizează
    scorul de material al oponentului după o singură mutare;
  \item \textbf{Hard} — evaluare la adâncimea 2: se simulează mutarea proprie și
    răspunsul optim al adversarului, alegând mutarea cu cel mai bun rezultat final.
\end{itemize}

Funcția de evaluare calculează suma valorilor pieselor rămase pe tablă
(pion=1, cal/nebun=3, tură=5, regină=9), din perspectiva culorii albe.

Dificultățile mai ridicate se deblochează automat pe baza victoriilor acumulate,
conform Tabelului~\ref{tab:chess_rewards}:

\begin{table}[H]
\centering
\caption{Deblocarea și recompensele șahului clasic}
\label{tab:chess_rewards}
\begin{tabular}{lcc}
\toprule
\textbf{Dificultate} & \textbf{Condiție deblocare} & \textbf{Recompensă victorie} \\
\midrule
Easy & Disponibil inițial & 12 monede \\
Normal & 3 victorii pe Easy & 24 monede + obiect inventar \\
Hard & 10 victorii pe Normal & 40 monede + obiect inventar \\
\bottomrule
\end{tabular}
\end{table}

\section{Puzzle-uri de șah tactice (mat în $N$ mutări)}
\label{sec:puzzles}

Pe lângă partidele complete, aplicația include un modul de \textbf{puzzle-uri tactice}
(\texttt{PuzzleScreen}), în care utilizatorul trebuie să găsească secvența de mat forțat
dintr-o poziție dată. Puzzle-urile sunt selecte din colecția publică \textit{wtharvey.com}
și acoperă trei niveluri de dificultate: \textbf{mat în 2}, \textbf{mat în 3} și
\textbf{mat în 4} mutări (Tabelul~\ref{tab:puzzle_rewards}). Un puzzle este prezentat
aleatoriu la fiecare deschidere a ecranului și la apăsarea butoanelor \textit{Next}/\textit{Skip}.

\begin{figure}[H]
\centering
\includegraphics[width=0.45\textwidth]{puzzle_screen.jpeg}
\caption{Ecranul de puzzle-uri de șah, cu o poziție de mat de rezolvat.\label{fig:puzzle}}
\end{figure}

\subsection{Pregătirea și validarea offline a puzzle-urilor}

Fișierele sursă conțin, pentru fiecare puzzle, o descriere, poziția în format FEN și soluția
în notație algebrică standard (SAN). Un script offline procesează aceste fișiere și produce
un fișier JSON compact (aproximativ 39\,KB pentru 70 de puzzle-uri) inclus în aplicație.
Pentru fiecare candidat, scriptul:
\begin{enumerate}
  \item normalizează câmpul FEN (numărul mutării este corectat la o valoare validă);
  \item rejoacă întreaga soluție cu \texttt{chess.js}, convertind fiecare mutare în format
    UCI;
  \item verifică faptul că secvența se termină efectiv în \textbf{șah-mat} și că adâncimea
    coincide cu cea declarată (mat în $N$);
  \item elimină puzzle-urile cu sub-promovări (promovare în altă piesă decât regina), pe care
    interfața curentă nu le poate exprima.
\end{enumerate}
Astfel, fiecare puzzle livrat este garantat \textbf{corect și solvabil}, validarea fiind
făcută complet \textit{offline}, fără un motor de șah extern.

\subsection{Mecanica de rezolvare}

Rezolvarea urmează \textbf{linia principală} a soluției validate: utilizatorul execută
mutarea așteptată, iar aplicația joacă automat răspunsul adversarului din soluție și continuă
până la mat, actualizând indicatorul ``mat în $N$''. O mutare legală dar greșită este
respinsă (cu \textit{feedback} sonor) și permite reîncercarea; o mutare care dă \emph{direct}
mat este de asemenea acceptată. Tabla se orientează automat după culoarea aflată la mutare,
iar butonul \textit{Hint} afișează mutarea corectă următoare. Recompensele cresc cu
dificultatea (Tabelul~\ref{tab:puzzle_rewards}), fiind acordate o singură dată per puzzle.

\begin{table}[H]
\centering
\caption{Recompensele puzzle-urilor de șah}
\label{tab:puzzle_rewards}
\begin{tabular}{lcc}
\toprule
\textbf{Dificultate} & \textbf{Număr puzzle-uri} & \textbf{Recompensă} \\
\midrule
Mat în 2 & 25 & 15 monede + 25 XP \\
Mat în 3 & 25 & 25 monede + 25 XP \\
Mat în 4 & 20 & 35 monede + 25 XP \\
\bottomrule
\end{tabular}
\end{table}

% ============================================================
%  CAPITOLUL 8 – TESTARE ȘI EVALUARE
% ============================================================
\chapter{Testare și evaluare}

\section{Strategia de testare}

Aplicația a fost testată la multiple niveluri:
\begin{enumerate}
  \item \textbf{Testare unitară} — un test de bază de randare a componentei rădăcină,
    scris cu Jest și React Test Renderer;
  \item \textbf{Testare manuală funcțională} — verificarea tuturor fluxurilor utilizator
    pe dispozitive Android fizice;
  \item \textbf{Verificări de integritate în cod} — validări explicite la rulare (ex.
    numărul de mutări din dicționarul Maia, starea modelelor ONNX).
\end{enumerate}

\section{Testare unitară}

Proiectul include teste unitare de bază în directorul \texttt{\_\_tests\_\_/App.test.tsx}
folosind Jest (\texttt{@types/jest 29.5.13}) și \texttt{react-test-renderer}. Configurația
Jest este definită în \texttt{jest.config.js} cu preset-ul
\texttt{react-native}.

\section{Testarea inferenței ONNX}

Ambele modele ONNX sunt încărcate asincron la montarea componentei corespunzătoare,
prin \texttt{InferenceSession.create()}. Indicatorii vizuali de stare
(\texttt{isModelLoading}, \texttt{isModelReady}, \texttt{modelsError}) sunt expuși
prin hook-urile \texttt{useOnnxClassifier} și \texttt{useMaiaChessEngine}, permițând
UI-ului să afișeze un spinner de încărcare și să dezactiveze interacțiunile cât timp
modelul nu este pregătit.

O verificare de integritate explicită este prezentă în codul sursă pentru modelul Maia:
dacă fișierul \texttt{all\_moves.json} nu conține exact 1880 de intrări (numărul oficial
din specificația Maia), se emite un mesaj de eroare critic în consolă. Aceasta asigură
că dicționarul de mutări utilizat este cel corect din repository-ul oficial Maia.

\section{Comportamentul hook-ului de geolocalizare}

Hook-ul \texttt{useGeolocation} (\texttt{src/hooks/useGeolocation.ts}) implementează
obținerea locației în mai mulți pași, cu un mecanism de rezervă explicit:

\begin{enumerate}
  \item \textbf{Solicitare permisiune} \textemdash{} pe Android se afișează un
    dialog de sistem cu titlul \textit{``Location Permission''} care explică scopul
    utilizării GPS-ului; pe iOS se apelează \texttt{Geolocation.requestAuthorization()};
  \item \textbf{Obținere coordonate GPS} \textemdash{} se apelează
    \texttt{Geolocation.getCurrentPosition()} cu parametrii:
    \texttt{enableHighAccuracy: false}, \texttt{timeout: 15000 ms},
    \texttt{maximumAge: 10000 ms};
  \item \textbf{Reverse geocoding Nominatim} \textemdash{} coordonatele sunt
    trimise la \texttt{nominatim.openstreetmap.org/reverse}; din câmpul
    \texttt{address} al răspunsului JSON se extrage denumirea localității prin
    ierarhia: \texttt{city} $\rightarrow$ \texttt{town} $\rightarrow$
    \texttt{village} $\rightarrow$ \texttt{county};
  \item \textbf{Rezervă la eșec} \textemdash{} un bloc \texttt{catch} explicit
    salvează locația cu doar latitudinea și longitudinea, fără denumire,
    asigurând că funcționalitatea nu se blochează în absența conectivității.
\end{enumerate}

Permisiunea \texttt{ACCESS\_FINE\_LOCATION} este solicitată la montarea ecranului
principal (\texttt{HomeScreen}). Sesiunile afișează fie denumirea localității (dacă
geocodarea a reușit), fie coordonatele numerice în formatul
\texttt{Lat: X.XXXX, Lon: Y.YYYY}.

\section{Mecanismul de persistență a datelor}

Persistența cu AsyncStorage este proiectată să fie robustă prin două
mecanisme complementare vizibile în codul sursă:
\begin{itemize}
  \item \textbf{Scriere la fiecare modificare de stare} \textemdash{} un \texttt{useEffect}
    care depend de \texttt{state} apelează asincron \texttt{AsyncStorage.setItem}
    la fiecare acțiune de reducer;
  \item \textbf{Versioning} \textemdash{} cheia de stocare \texttt{retro-arcade-study-state-v1}
    este versionată, permițând migrări viitoare de schemă fără pierderea datelor
    existente. La restaurare (acțiunea \texttt{HYDRATE}), starea salvată este
    îmbinată cu \texttt{initialState} ca baseline, protejând împotriva câmpurilor
    lipsă apărute în versiuni mai noi ale aplicației.
\end{itemize}

\section{Testarea jocului de șah Maia}

Corectitudinea integrării Maia Chess a fost validată prin:
\begin{itemize}
  \item Verificarea că toate mutările executate de AI sunt mutări legale conform
    \texttt{chess.js};
  \item Verificarea tehnicii de oglindire FEN prin testarea pozițiilor specifice cu
    negrul la mutare;
  \item Verificarea că numărul de intrări în \texttt{all\_moves.json} este exact 1880
    (conform specificației Maia) — validare prezentă explicit în codul sursă la
    inițializarea modulului.
\end{itemize}

\section{Validarea offline a datelor de joc}

O parte din conținutul aplicației este validat \textbf{algoritmic, în avans}, în loc să fie
verificat manual, ceea ce crește încrederea în corectitudine:
\begin{itemize}
  \item \textbf{Puzzle-urile de șah} — fiecare poziție și soluție din colecția importată este
    rejucată cu \texttt{chess.js}, iar puzzle-ul este păstrat doar dacă secvența se termină în
    șah-mat la adâncimea declarată (Secțiunea~\ref{sec:puzzles}); puzzle-urile invalide sau cu
    sub-promovări sunt eliminate automat.
  \item \textbf{Cartea de deschideri de alb} — pentru fiecare deschidere, legalitatea celei
    de-a doua mutări a fost verificată împotriva unui set variat de răspunsuri ale negrului,
    garantând că AI-ul nu ajunge niciodată într-o mutare ilegală în deschidere.
  \item \textbf{Conținutul de învățare} — la importul exercițiilor au fost verificate
    integritatea câmpurilor, potrivirea răspunsului corect cu una dintre variante și unicitatea
    identificatorilor.
\end{itemize}
Aceste verificări sunt rulate la pregătirea datelor (\textit{build-time}), astfel încât
aplicația livrată conține doar conținut validat.

% ============================================================
%  CAPITOLUL 9 – CONCLUZII ȘI DIRECȚII VIITOARE
% ============================================================
\chapter{Concluzii și direcții viitoare}

\section{Concluzii}

Prezenta lucrare a descris proiectarea și implementarea completă a aplicației mobile
\textbf{FocusQuest}, demonstrând viabilitatea unui produs software complex care combină
multiple domenii: productivitate, gamificare, educație și inteligență artificială
\textit{on-device}.

Principalele realizări tehnice ale proiectului sunt:

\begin{enumerate}
  \item \textbf{Integrarea reușită a ONNX Runtime în React Native} — ambele modele
    (ResNet-50 și Maia Chess) funcționează complet offline pe dispozitive Android,
    cu timpi de răspuns acceptabili pentru interacțiunea utilizatorului.

  \item \textbf{Implementarea corectă a motorului Maia Chess} — inclusiv rezolvarea
    problemelor critice de oglindire FEN, codificarea celor 18 canale de intrare,
    maparea Elo la categorii și selecția mutărilor prin softmax pe mutările legale.

  \item \textbf{Un sistem de gamificare complet și coerent} — cu economie virtuală
    (monede), experiență (XP) și niveluri progresive, quest-uri zilnice, efecte sonore,
    deblocarea progresivă a conținutului, 42 de realizări, inventar cu 78 de
    obiecte și animații RPG.

  \item \textbf{Un modul educațional funcțional} — cu 673 de exerciții matematice reale,
    un mod de examen simulat, randare matematică cu KaTeX şi sistem de progres persistent.

  \item \textbf{Un modul de puzzle-uri de șah tactice} — cu 70 de probleme de tip mat în
    2--4 mutări validate algoritmic și un mecanism de rezolvare pe linia principală.

  \item \textbf{O arhitectură solidă} — bazată pe React Context + useReducer cu
    persistență AsyncStorage, TypeScript strict și navigare tipizată.
\end{enumerate}

Aplicația demonstrează că un singur dezvoltator poate construi, în cadrul unui proiect
de licență, o aplicație mobilă funcțională și complexă care integrează modele de
machine learning \textit{on-device}, fără a depinde de infrastructură cloud.

\section{Limitări curente}

În forma sa actuală, aplicația prezintă câteva limitări asumate, identificate pe parcursul
dezvoltării și testării:
\begin{itemize}
  \item \textbf{Dependența de internet pentru randarea matematică} — biblioteca KaTeX este
    încărcată din CDN, ceea ce face ca modulul de învățare să necesite o conexiune activă
    pentru afișarea corectă a formulelor; restul aplicației funcționează \textit{offline}.
  \item \textbf{Suport audio limitat la Android} — muzica de fundal este redată din resursele
    native Android (\texttt{res/raw}); portarea completă pe iOS necesită includerea
    fișierelor în bundle-ul corespunzător.
  \item \textbf{ELO fix al jucătorului în Maia Chess} — parametrul \texttt{elo\_oppo} este
    fixat la 1600, indiferent de nivelul real al utilizatorului, ceea ce menține constant
    stilul adaptiv al adversarului AI.
  \item \textbf{Acoperirea parțială a testării automate} — validarea s-a bazat preponderent
    pe testare manuală pe dispozitiv fizic și emulator, completată de un test unitar de
    randare și de verificări de integritate în cod.
\end{itemize}
Aceste limitări nu afectează funcționalitatea de bază a aplicației și constituie puncte de
plecare naturale pentru direcțiile viitoare de dezvoltare prezentate în continuare.

\section{Direcții viitoare de dezvoltare}

\subsection{Extinderea conținutului educațional}
Adăugarea de noi discipline (fizică mecanică, chimie organică, informatică) cu seturi
suplimentare de exerciții matematice şi ştiințifice.

\subsection{Repetiție spațiată (spaced repetition)}
Introducerea unui algoritm de repetiție spațiată (de exemplu SM-2 sau sistemul Leitner)
pentru a reprograma automat exercițiile și intrările de jurnal greșite sau dificile la
intervale crescătoare, valorificând câmpul \texttt{lastReviewedAt} deja prezent în model
și maximizând retenția pe termen lung.

\subsection{Rating adaptiv al jucătorului}
Înlocuirea parametrului fix \texttt{elo\_oppo = 1600} cu un rating al jucătorului calculat
adaptiv (de exemplu prin sistemul Glicko-2) pe baza rezultatelor împotriva nivelurilor Maia,
oferind atât o adversitate mai realistă, cât și un indicator de progres suplimentar.

\subsection{Sincronizare cloud opțională}
Implementarea unui backend opțional (Firebase / Supabase) pentru sincronizarea
progresului pe mai multe dispozitive și funcționalitate de leaderboard global.

\subsection{Modul multiplayer pentru șah}
Adăugarea posibilității de a juca șah contra altor utilizatori ai aplicației în timp
real, folosind WebSocket-uri.

\subsection{Îmbunătățirea clasificatorului AI}
Fine-tuning al modelului ResNet pe un set de date mai relevant contextului educațional
(ex. recunoașterea tipurilor de probleme matematice din fotografii de caiete).

\subsection{Notificări push}
Implementarea notificărilor zilnice de reamintire a streak-ului de concentrare și a
exercițiilor de rezolvat.

\subsection{Suport iOS complet}
Testarea și corectarea comportamentului pe dispozitive iOS, cu publicarea în App Store.

\subsection{Accesibilitate}
Adăugarea suportului pentru modul de contrast ridicat și compatibilitate cu cititorii
de ecran (TalkBack / VoiceOver).

% ============================================================
%  BIBLIOGRAFIE
% ============================================================
\begin{thebibliography}{99}
\addcontentsline{toc}{chapter}{Bibliografie}

\bibitem{reactnative}
Meta Platforms, Inc.,
\textit{React Native -- Build native apps using React},
\url{https://reactnative.dev}, 2024.

\bibitem{onnxruntime}
Microsoft Corporation,
\textit{ONNX Runtime: cross-platform, high performance ML inferencing and training accelerator},
\url{https://onnxruntime.ai}, 2024.

\bibitem{onnxmodelzoo}
ONNX Community,
\textit{ONNX Model Zoo -- ResNet50-v1-7},
\url{https://github.com/onnx/models}, 2024.

\bibitem{maiacChess}
R.~McIlroy-Young, S.~Sen, J.~Kleinberg, A.~Anderson,
\textit{Aligning Superhuman AI with Human Behavior: Chess as a Model System},
Proceedings of the 26th ACM SIGKDD International Conference on Knowledge
Discovery \& Data Mining, 2020.

\bibitem{maiaGitHub}
Maia Chess Project,
\textit{Maia Chess -- Human-Like AI for Chess},
\url{https://github.com/CSSLab/maia-chess}, 2020.

\bibitem{resnet}
K.~He, X.~Zhang, S.~Ren, J.~Sun,
\textit{Deep Residual Learning for Image Recognition},
Proceedings of the IEEE Conference on Computer Vision and Pattern Recognition
(CVPR), pp. 770–778, 2016.

\bibitem{reactnavigation}
React Navigation Contributors,
\textit{React Navigation -- Routing and navigation for Expo and React Native apps},
\url{https://reactnavigation.org}, 2024.

\bibitem{asyncstorage}
React Native Community,
\textit{@react-native-async-storage/async-storage},
\url{https://react-native-async-storage.github.io/async-storage/}, 2024.

\bibitem{chessjs}
J.~Huba,
\textit{chess.js -- A JavaScript chess library},
\url{https://github.com/jhlywa/chess.js}, 2024.

\bibitem{katex}
Khan Academy,
\textit{KaTeX -- The fastest math typesetting library for the web},
\url{https://katex.org}, 2024.

\bibitem{habitica}
HabitRPG, Inc.,
\textit{Habitica -- Gamify Your Life},
\url{https://habitica.com}, 2024.

\bibitem{forest}
Seekrtech Co. Ltd.,
\textit{Forest -- Stay Focused},
\url{https://www.forestapp.cc}, 2024.

\bibitem{lichess}
Lichess.org,
\textit{Lichess -- Free, Open Source Chess Server},
\url{https://lichess.org}, 2024.

\bibitem{stockfish}
Stockfish Team,
\textit{Stockfish -- Open Source Chess Engine},
\url{https://stockfishchess.org}, 2024.

\bibitem{duolingo}
Duolingo, Inc.,
\textit{Duolingo -- Learn a language for free},
\url{https://www.duolingo.com}, 2024.

\bibitem{nominatim}
OpenStreetMap Contributors,
\textit{Nominatim -- OpenStreetMap Geocoding},
\url{https://nominatim.openstreetmap.org}, 2024.

\bibitem{typescript}
Microsoft Corporation,
\textit{TypeScript -- JavaScript With Syntax For Types},
\url{https://www.typescriptlang.org}, 2024.

\bibitem{imageNet}
J.~Deng, W.~Dong, R.~Socher, L.~Li, K.~Li, L.~Fei-Fei,
\textit{ImageNet: A Large-Scale Hierarchical Image Database},
Proceedings of CVPR, 2009.

\bibitem{onnxformat}
ONNX Community,
\textit{Open Neural Network Exchange (ONNX)},
\url{https://onnx.ai}, 2024.

\bibitem{minimax}
S.~Russell, P.~Norvig,
\textit{Artificial Intelligence: A Modern Approach}, 4th edition,
Pearson, 2020.

\bibitem{sudoku}
A.~Maji, S.~Pal,
\textit{Sudoku Solving Algorithms: A Brief Review},
International Journal of Computer Science and Engineering, 2013.

\end{thebibliography}


\end{document}
